% The next line makes it so it compiles the main file if I hit the compile button
% when editing this file in TeXworks.
% !TEX root = ../AIDMA.tex
%%%%%%%%%%%%%%%%%%%%%%%%%%%%%%%%%%%%

\chapter{Sequences, Summations, and Matrices}\label{ch:seqsummat}

\section{Sequences}\label{section:seq}
\begin{df}\index{sequence}
A {\bf sequence} of real numbers is a function whose domain is the set of
natural numbers and whose output is a subset of the real numbers. We
usually denote a sequence by one of the notations
$$a_0, a_1, a_2, \ldots 
$$ 
or
$$ \{a_n\}_{n=0} ^{+\infty}  $$
or
$$ \{a_n\}.$$
The last notation is just a shorthand for the second notation.
\end{df}

\begin{rem}
Since sequences are functions, sometimes function notation is used.  That is, $a(n)$
instead of $a_n$.
\end{rem}

We will be mostly interested in two types of sequences. 
The first type are sequences that have an explicit formula for their $n$-th term.
They are said to be in {\em closed form}.

\begin{exa}
Let $a_n = 1 - \frac{1}{2^n}, n = 0, 1, \ldots$. Then $\{a_n\}_{n =
0} ^{+\infty}$ is a sequence for which we have an explicit formula
for the $n$-th term. The first five terms are
$$
\begin{array}{lllllll} 
a_0 & = & 1 -\frac{1}{2^0} &=&1-1 & = & 0, \\[1ex]
 a_1 & = & 1 -\frac{1}{2^1}&=&1-\frac{1}{2} &= & \frac{1}{2}, \\[1ex]
a_2 & = & 1 -\frac{1}{2^2}&=&1-\frac{1}{4} &= & \frac{3}{4}, \\[1ex]
a_3 & = & 1 -\frac{1}{2^3}&=&1-\frac{1}{8} &= & \frac{7}{8}, \\[1ex]
a_4 & = & 1 -\frac{1}{2^4}&=&1-\frac{1}{16} &= & \frac{15}{16}. \\
\end{array}
$$
\end{exa}

\begin{rem}
Sometimes we may not start at $n = 0$. In that case we may write
$$a_m, a_{m+1}, a_{m+2}, \ldots,$$ 
or
$$ \{a_n\}_{n=m} ^{+\infty}\ ,$$
where $m$ is a non-negative integer.  
Most sequences we will deal with will start with $m=0$ or $m=1$.
\end{rem}

\begin{exer}
Let $\{x_n\}$ be the sequence defined by $x_n = 1 + (-2)^n, n = 0, 1, 2, \ldots$.
Find the first five terms of $\{x_n\}$.
\begin{lenum}
\item $x_0=\blankline{2}$
\item $x_1=\blankline{2}$
\item $x_2=\blankline{2}$
\item $x_3=\blankline{2}$
\item $x_4=\blankline{2}$
\end{lenum}
\begin{answer}
 (a)$x_0=1+(-2)^0 = 1 + 1 = 2$
 (b)$x_1=1+(-2)^1 = 1-2=-1$
 (c)$x_2=1+(-2)^2=1+4= 5$
 (d)$x_3=1+(-2)^3=1-8=-7$
 (e)$x_4=1+(-2)^4=1+16=17$
\end{answer}
\end{exer}

\begin{exer}
Find the first five terms of the following sequences.
\begin{lenum}
  \setlength{\itemsep}{.35in}
\item 
${\displaystyle x_n = 1 + \left(-\frac{1}{2}\right)^n, n = 0, 1, 2, \ldots}$

\noindent
$x_0=\blankline{1.5}$\, \,  $x_1=\blankline{1.5}$\, \,  $x_2=\blankline{1.5}$

\noindent
$x_3=\blankline{1.5}$\, \,  $x_4=\blankline{1.5}$
\item 
${\displaystyle x_n = n! + 1, n = 0, 1, 2, \ldots}$

\noindent
$x_0=\blankline{1.5}$\, \,  $x_1=\blankline{1.5}$\, \,  $x_2=\blankline{1.5}$

\noindent
$x_3=\blankline{1.5}$\, \,  $x_4=\blankline{1.5}$ 
\item 
${\displaystyle x_n = \dfrac{1}{n! + (-1)^n}, n =  2, 3, 4,\ldots}$

\noindent
$x_2=\blankline{1.5}$\, \,  $x_3=\blankline{1.5}$\, \,  $x_4=\blankline{1.5}$

\noindent
$x_5=\blankline{1.5}$\, \,  $x_6=\blankline{1.5}$
\item
${\displaystyle x_n = \left(1 +\dfrac{1}{n}\right)^{n}, n = 1, 2,\ldots }$

\noindent
$x_1=\blankline{1.5}$\, \,  $x_2=\blankline{1.5}$\, \,  $x_3=\blankline{1.5}$

\noindent
$x_4=\blankline{1.5}$\, \,  $x_5=\blankline{1.5}$ 
\end{lenum}
\vspace*{.25in}
\begin{answer}
We will just provide the final answer for these.  If you can't get these answers, you may need to brush 
up on your algebra skills.
(a) $2$, $1/2$,$5/4$, $7/8$, $17/16$; 
(b) $2$, $2$, $3$, $7$, $25$; 
(c) $1/3$, $1/5$, $1/25$, $1/119$, $1/721$; 
(d) $2$, $9/4$, $64/27$, $625/256$, $7776/3125$
\end{answer}
\end{exer}


The second type of sequence are defined using {\em recurrence relations}.

\begin{df}
\index{recurrence relations!definition}
A {\bf recurrence relation} is an equation
that defines each term of a sequence based on one or more
previous terms of the sequence.
More specifically, a recurrence relation for a sequence $\{a_n\}$
will define $a_n$ based on (some of) the values of 
$a_0$, $a_1$, \ldots, $a_{n-1}$.
\end{df}

\begin{exa}\label{exa:rec1}
Let $\displaystyle x_0 = 1,\ \  x_n = \left(1 + \frac{1}{n}\right)x_{n-1},\mbox{ for } n = 1, 2, \ldots .$ 
Then $\{x_n\}_{n = 0} ^{+\infty}$ is a 
recursively defined sequence. The terms $x_1, x_2, \ldots , x_5$ are
$$\begin{array}{lllllllll}
x_1 & = & \left(1 +\frac{1}{1}\right)x_0 &=& \left(1 +\frac{1}{1}\right)1 &= &1+1&= & 2. \\[1ex]
x_2 & = & \left(1 +\frac{1}{2}\right)x_1 &=& \left(1 +\frac{1}{2}\right)2 &= &2+1&= & 3. \\[1ex]
x_3 & = & \left(1 +\frac{1}{3}\right)x_2 &=& \left(1 +\frac{1}{3}\right)3 &= &3+1&= & 4. \\[1ex]
x_4 & = & \left(1 +\frac{1}{4}\right)x_3 &=& \left(1 +\frac{1}{4}\right)4 &= &4+1&= & 5. \\[1ex]
x_5 & = & \left(1 +\frac{1}{5}\right)x_4 &=& \left(1 +\frac{1}{5}\right)5 &= &5+1&= & 6. \\[1ex]
\end{array}
$$
\end{exa}


Notice that in the previous example, we gave an explicit definition of $x_0$.
This is called an {\em initial condition}.  In order to specify a 
sequence, a recurrence relation needs one or more initial conditions. 
Without them, we have an abstract definition of a sequence,
but cannot compute any values since there is no ``starting point.''
Also note that different initial conditions can be specified for the same
recurrence relation, resulting in different sequences being generated.

When we find an explicit formula (or {\em closed formula}) for a recurrence relation, 
we say we have {\em solved} the recurrence relation.

\begin{exa}
Given the values we computed in Example~\ref{exa:rec1}, 
it seems relatively clear that $x_n=n+1$ is a solution for that recurrence relation. 
\end{exa}

\begin{rem}
It is important to be careful about jumping to conclusions too quickly when solving
recurrence relations.\footnote{These comments also apply to other problems that involve seeing a
pattern and finding an explicit formula.}
Although it turns out that in the previous example, $x_n=n+1$ is the correct closed form
(we will prove it shortly),
just because it works for the first 5 terms does not necessarily imply that the pattern continues.
\end{rem}

\begin{exer}\label{exer:5pow}
Let $\{x_n\}$ be the sequence defined by 
$$x_0=1, x_n=5\cdot x_{n-1}, \mbox{ for }  n=1,2,\ldots.$$
Find a closed form for $x_n$.  (Hint: Start by computing $x_1$, $x_2$, $x_3$, etc. until you see the pattern.)

\vspace*{1.9in}
\begin{answer}
Notice that $x_0=1$, $x_1=5\cdot 1 = 5$, $x_2=5\cdot 5=5^2$, $x_3=5\cdot 5^2 = 5^3$, etc.
Looking back, we can see that $1=5^0$, so $x_0=5^0$.  Also, $x_1=5=5^1$.  So it seems likely that the
solution is $x_n=5^n$.  This is {\em not} a proof, though!
\end{answer}
\end{exer}

\begin{exer}\label{exer:factrec}
Let $\{x_n\}$ be the sequence defined by 
$$x_0=1, x_n=n\cdot x_{n-1}, \mbox{ for }  n=1,2,\ldots.$$
Find a closed form for $x_n$.


\vspace*{1.9in}
\begin{answer}
Notice that $x_0=1$, 
$x_1=1\cdot 1 = 1$, 
$x_2=2\cdot 1=2$, 
$x_3=3\cdot 2 = 6$,
$x_4=4\cdot 6 = 24$,
$x_3=5\cdot 24 = 120$, etc.
Written this way, no obvious pattern is emerging.  Sometimes {\em how} you write the numbers matters.
Let's try this again:
$x_1=1\cdot 1 = 1!$, 
$x_2=2\cdot 1=2!$, 
$x_3=3\cdot 2\cdot 1=3!$,
$x_4=4\cdot 3\cdot 2\cdot 1=4!$,
$x_3=5\cdot 4\cdot 3\cdot 2\cdot 1=5!$, etc.
Now we can see that $x_n=n!$ is a likely solution.  Again, this isn't a proof.
\end{answer}
\end{exer}

\begin{eval}
Define $\{a_n\}$ by $a(0)=1$, $a(1)=2$, and 
$$a_n=\left\lfloor \frac{1+\sqrt{5}}{2}\times a_{n-1}\right\rfloor+a_{n-2}$$
for $n\geq 2$.  Find a closed form for $a_n$.
\begin{ssol}
We can see that 
$$\begin{array}{lllllll} 
a_2 & = & \left\lfloor \frac{1+\sqrt{5}}{2}\times a_1\right\rfloor+a_0&=
& \left\lfloor \frac{1+\sqrt{5}}{2}\times 2\right\rfloor+1&=&4\\[1ex]
a_3 & = &\left\lfloor \frac{1+\sqrt{5}}{2}\times a_2\right\rfloor+a_1&=
& \left\lfloor \frac{1+\sqrt{5}}{2}\times 4\right\rfloor+2&=&8 \\[1ex]
a_4 & = &\left\lfloor \frac{1+\sqrt{5}}{2}\times a_3\right\rfloor+a_2&=
& \left\lfloor \frac{1+\sqrt{5}}{2}\times 8\right\rfloor+4&=&16
\end{array}
$$
(You can verify these with a calculator).
At this point it seems relatively clear that $a_n=2^n$.  
\end{ssol}
\evallinetwo
\begin{answer}
Their calculations are correct (Did you check them with a calculator?  You {\em should} have!
How else can you tell whether or not their solution is correct?).
So it does seem like $a_n=2^n$ is the correct solution.
However, 
$$\begin{array}{lllllll} 
a_5  = \left\lfloor \frac{1+\sqrt{5}}{2}\times a_4\right\rfloor+a_3=
 \left\lfloor \frac{1+\sqrt{5}}{2}\times 16\right\rfloor+8=33\neq 2^5
\end{array}$$
so the solution that seems `obvious' turns out to incorrect.
We won't give the actual solution since the point of this example is to demonstrate
that just because a pattern holds for the first several terms of a sequence, it does not
guarantee that it holds for the whole sequence.
\end{answer}
\end{eval}

Did you catch what happened in the previous Evaluate exercise?  The `obvious' solution
wasn't correct.  If you missed this, go back and read the solution. 

Generally speaking, you need to {\em prove} that the closed form is correct.
One way to do this is to plug it back into the recursive definition.
If we can plug it into the right hand side of the recursive definition and are able to simplify it to
the left hand side, then it must be a solution.  We also have to verify that it works for the
initial condition(s).

As an analogy, how do you know that $x=-1$ is a solution to the equation $x^2+2x+1=0$?
You plug it in to get $(-1)^2+2(-1)+1=1-2+1=0$.  Since we got $0$, $x=-1$ is a solution.
We do something similar for recurrence relations, 
except that what we are plugging in is a formula instead of just a number.

\begin{exa}\label{exa:firstVerify}
Prove that $x_n=n+1$ is a solution to the recurrence relation given by 
$$x_0 = 1,\ \  x_n = \left(1 + \frac{1}{n}\right)x_{n-1},\ \ \ n
= 1, 2, \ldots .$$
\begin{pf}
To prove that $x_n=n+1$ is a solution for $n\geq 0$, 
we need to show two things.  First, that it works for the initial condition.
Since $x_0=1=0+1$, it works for the initial condition.
Second, that if we plug it into the right hand side of the recursive definition, that
we can simplify it to $x_n$.  Doing so, we get 
\begin{eqnarray*}
\left(1 + \frac{1}{n}\right)x_{n-1} &= &\left(1 + \frac{1}{n}\right)((n-1)+1)\\
&=&\left(\frac{n+1}{n}\right)n\\
&=&n+1\\
&=&x_n
\end{eqnarray*}
Since plugging the solution back in verifies the recurrence relation, 
$x_n=n+1$ is a solution to the recurrence relation.

If you are confused by the first step of algebra, remember that we are assuming that $x_n=n+1$ for $n\geq 0$.
Thus, $x_{n-1}=(n-1)+1=n$, since we are just plugging in $n-1$ instead of $n$.
\end{pf}
\end{exa}

\begin{exer}
Prove that your solution to Exercise~\ref{exer:5pow} is correct.

\vspace*{2.5in}
\begin{answer}
Hopefully you came up with the solution $x_n=5^n$.  
Since $x_0=1=5^0$, it works for the initial condition.
If we plug this back into the
right hand side of $x_n=5\cdot x_{n-1}$,  we get
\begin{eqnarray*}
5\cdot x_{n-1}&=&5\cdot 5^{n-1}\\
&=&5^n\\
&=&x_n,
\end{eqnarray*}
which verifies the formula.  Therefore $x_n=5^n$ is the solution.
\end{answer}
\end{exer}

\begin{exer}
Prove that your solution to Exercise~\ref{exer:factrec} is correct.

\vspace*{2.5in}
\begin{answer}
Hopefully you came up with the solution $x_n=n!$.  
Since $x_0=1=0!$, it works for the initial condition.
If we plug this back into the
right hand side of $x_n=n\cdot x_{n-1}$,  we get
\begin{eqnarray*}
n\cdot x_{n-1}&=&n\cdot (n-1)!\\
&=&n!\\
&=&x_n,
\end{eqnarray*}
which verifies the formula.  Therefore $x_n=n!$ is the solution.
\end{answer}
\end{exer}

\newpage
\begin{eval}\label{eval:ferzle}
Determine what \verb+ferzle(n)+ (below) returns for $n=0, 1, 2, 3, 4$
and then re-write \verb+ferzle+ without using recursion, making it 
as efficient as possible.\footnote{Although we have
not formally covered recursion yet, we expect that you have seen it before and
know enough to follow this example. If not, ask your instructor or a friend for help.}
\begin{code}
	int ferzle(int n) {
	    if(n<=0) {
	        return 3;
	    } else {
	        return ferzle(n-1) + 2;
	    }
	}
\end{code}
\begin{ssol}
First, we can see that ferzle(0) returns 3 since it executes the code in the if statement.
ferzle(1) returns ferzle(0)+2, which is $3+2=5$.
ferzle(2) returns ferzle(1)+2, which is $5+2=7$.
ferzle(3) returns ferzle(2)+2, which is $7+2=9$.
ferzle(4) returns ferzle(3)+2, which is $9+2=11$.
Notice that $11=2*4+3$, $9=2*3+3$, $7=2*2+3$, $5=2*1+3$, and $3=2*0+3$.
From this, it is pretty clear that ferzle(n) returns $2n+3$.
Thus, my simplified function is as follows:
\begin{verbatim}
    int ferzle(int n) {
        return 2*n+3;
    }
\end{verbatim}
\end{ssol}
\evallinethree
\begin{answer}
The computations are correct, the conclusion is correct, but unfortunately,
the final code has a serious problem.  It works {\em most} of the time,
but it does not deal with negative values correctly.  
It should return $3$ for all negative
values, but it continues to use the formula.
The problem is they forgot to even consider what the function does for negative values of $n$.
They probably could have formatted their answer better, too.  It's difficult to follow in paragraph
form.  They could have put the various values of \verb+ferzle(n)+ each on their own line
and presented it mathematically instead of in sentence form.
For instance, instead of `ferzle(1) returns ferzle(0)+2, which is $3+2=5$,'
they should have `$ferzle(1)=ferzle(0)+2=3+2=5$.'
It would have made it much easier to see the pattern.
\end{answer}
\end{eval}

\begin{exer}
Fix the code from the solution given in Evaluate~\ref{eval:ferzle} 
so that it still uses the closed form, but works correctly for all values of $n$.
\begin{code}
	int ferzle(int n) {
	 
	 
	 
	 
	 
	 
	 
	 
	}
\end{code}
\begin{answer}\begin{code}
	int ferzle(int n) {
		if(n<=0) {
		    return 3;
	    } else {
		    return 2*n+3;
		}
	}
\end{code}
\end{answer}
\end{exer}



A more complete discussion of solving recurrences appears in Chapter \ref{ch:recurrences}.

The following is a famous example of a recursively defined sequence that we will revisit several times.
\begin{exa}\index{Fibonacci numbers}\index{Fibonacci sequence}
The {\em Fibonacci sequence} is a sequence of numbers that is of interest
in various mathematical and computing applications.  
They are defined using the following recurrence relation:\footnote{In the remainder of the book, 
when you see $f_k$, you should assume it refers to the $k$-th Fibonacci number
unless otherwise specified.}
$$f_n=\left\{
\begin{array}{ll}
0&\mbox{ if $n$=0}\\
1&\mbox{ if $n$=1}\\
f_{n-1}+f_{n-2}&\mbox{ if }n>1\\
\end{array}
\right.
$$
In words, each Fibonacci number (beyond the first two) is the sum of the previous two.
The first few are $f_0=0$, $f_1=1$, 
\begin{eqnarray*}
  f_2 &=& f_1+f_0 = 1 + 0 =  1,\\
  f_3 &=&f_2+f_1  = 1 + 1 =  2,\\
  f_4 &=& f_3+f_2  = 2 + 1 =  3,\\
  f_5 &=& f_4+f_3  = 3 + 2 =  5,\\
  f_6 &=& f_5+f_4  = 5 + 3 =  8,\\
  f_7 &=& f_6+f_5  = 8 + 5 = 13.
\end{eqnarray*}
Later we will see the closed form for the Fibonacci sequence.
If you are really adventurous, you might consider trying to determine it yourself.
But be warned:  It is not a simple formula that you will come up with by just looking
at some of the Fibonacci numbers.
\end{exa}


\begin{df}
A sequence $\{a_n\} _{n=0} ^{+\infty}$ is said to 
be 
\begin{itemize}
\item
\index{increasing sequence}
{\bf increasing} if $a_{n}\leq a_{n + 1} \ \forall
n\in\naturals$
\item  
\index{strictly increasing sequence}
{\bf strictly increasing} if $a_n < a_{n + 1}
\ \forall n\in\naturals$
\item
\index{decreasing sequence}
{\bf decreasing} if $a_{n}\geq a_{n + 1} \ \forall
n\in\naturals$
\item
\index{strictly decreasing sequence}
{\bf strictly decreasing} if $a_n > a_{n + 1}
\ \forall n\in\naturals$
\end{itemize}
Some people call these sequences 
{\bf non-decreasing}, {\bf increasing}, {\bf non-increasing}, and {\bf decreasing}, respectively. 

\index{monotonic sequence}\index{non-monotonic sequence}
A sequence is called {\bf monotonic} if it is any of these, 
and {\bf non-monotonic} if it is none of these.
\end{df}

\begin{exa}
Recall that $0! = 1$, $1! = 1$, $2! = 1\cdot 2 = 2$, $3! = 1\cdot 2
\cdot 3 = 6$, etc. Prove that the sequence $x_n = n!, n = 0, 1, 2,
\ldots$ is strictly increasing for $n \geq 1$.
 \begin{pf}   For $n > 1$ we have $$x_n = n! = n(n - 1)! = nx_{n - 1}
> x_{n - 1},
$$since $n>1$. This proves that the sequence is strictly
increasing.
\end{pf}
\end{exa}

\begin{ques}
Notice in this first example we concluded that the sequence is strictly increasing since we showed that
$x_n>x_{n-1}$.  But according to the definition we need to show that $x_n<x_{n+1}$.  
So did we do something wrong?  Explain.

\noindent
\answerlinethree
\begin{answer}
We didn't do anything wrong.  We wrote the inequality in the other order, and the indexes are one lower
than those given in the definition.  But that's O.K.  The definition is simply trying to convey the idea
that every term is strictly greater than the previous.  That is what we showed.  
We can show that  
$x_n<x_{n+1}$, $x_{n+1}>x_n$, $x_{n-1}<x_{n}$, or $x_n>x_{n-1}$.  They all mean essentially the same thing. 
The only difference is the order in which the inequalities are written (the first two and last two are
saying exactly the same thing--we just flipped the inequality) and what values of $n$ are valid.
For instance, if the sequence starts at $0$, then we need to assume
$n\geq 0$ for the first pair of inequalities and  $n\geq 1$ for the second pair. 
\end{answer}
\end{ques} 

\begin{exa}
Prove that the sequence $x_n = 2 + \dfrac{1}{2^n}$, $n =0, 1, 2,
\ldots$ is strictly decreasing.
 \begin{pf}   We have $$\begin{array}{lll}x_{n + 1} - x_n & = & \left(2
+ \dfrac{1}{2^{n+1}} \right) - \left(2 +
\dfrac{1}{2^n}\right) \\
& = & \dfrac{1}{2^{n+1}} -\dfrac{1}{2^n} \\
& = & -\dfrac{1}{2^{n+1}} \\
& < & 0.
\end{array}$$
Thus, $x_{n + 1} - x_n < 0$, so $x_n > x_n{+1},$
i.e., the sequence is strictly decreasing.
\end{pf}
\end{exa}

\begin{exer}
Prove that the sequence $x_n = \dfrac{n^2 + 1}{n}$, $n = 1, 2,
\ldots$ is strictly increasing.

\vspace*{3.1in}
\begin{answer}
If you got stuck on this one, first realize that $x_n=\dfrac{n^2 + 1}{n} = n + \dfrac{1}{n}$.
This form might make the algebra a little easier.  
Then, follow the technique of the previous example--show that
$x_{n + 1} - x_n >0$.  So, if necessary, go back and try again.
If you already attempted a proof, you may proceed to read the solution.

Notice that, $$ \begin{array}{lll}x_{n + 1} - x_n & = & \left(n + 1 +
\dfrac{1}{n + 1}\right) -\left(n  + \dfrac{1}{n
}\right) \\
& = & 1 + \dfrac{1}{n + 1} - \dfrac{1}{n} \\
& = & 1 -\dfrac{1}{n(n +1)} \\
& > & 0,\\
\end{array}$$
the last step since  $1/n(n+1)<1$ when $n\geq 1$.
Therefore, $x_{n + 1} - x_n > 0$, so $x_{n + 1}> x_n,$
i.e., the sequence is strictly increasing.
If your solution is significantly different than this, make sure you determine one way or another
if it is correct.
\end{answer}
\end{exer}


\begin{exer}
Decide whether the following sequences are increasing, strictly increasing,
decreasing, strictly decreasing, or non-monotonic.  
You do not need to prove your answer, but give a brief justification.
\begin{lenum}
\item $x_n = n,$ $n = 0, 1, 2, \ldots$

\noindent \answerlinetwo
\item $x_n = (-1)^nn,$ $n = 0, 1, 2, \ldots$

\noindent \answerlinetwo
\item $x_n = \dfrac{1}{n!},$ $n = 0, 1, 2, \ldots$

\noindent  \answerlinetwo
\item $x_n = \dfrac{n}{n + 1},$ $n = 0, 1, 2, \ldots$

\noindent \answerlinetwo
\item $x_n = n^2 - n,$ $n = 1, 2, \ldots$

\noindent \answerlinetwo
\item $x_n = n^2 - n,$ $n = 0, 1, 2, \ldots$

\noindent \answerlinetwo
\item $x_n = (-1)^n,$ $n = 0, 1, 2, \ldots$

\noindent \answerlinetwo
\item $x_n = 1 - \dfrac{1}{2^n},$ $n = 0, 1, 2, \ldots$

\noindent  \answerlinetwo
\item $x_n = 1 + \dfrac{1}{2^n},$ $n = 0, 1, 2, \ldots$

\noindent \answerlinetwo
\end{lenum}
\begin{answer}
We could go into much more detail than we do here, and hopefully you did when you wrote down
your solutions.  But we'll settle for short, informal arguments this time.
 (a) This is just a linear function.  It is {\bf strictly increasing}.
 (b) Since this keeps going from positive to negative to positive, etc. it is {\bf non-monotonic}.
 (c) We know that $n!$ is strictly increasing.  Since this is the reciprocal of that function, it
 is almost {\em strictly decreasing} (since we are dividing by a number that is getting larger). 
 However, since $1/0!=1/1!=1$, it is just {\bf decreasing}.
(d) This is getting closer to $1$ as $n$ increases.  It is {\bf strictly increasing} 
(e) This is $n(n-1)$.  $x_1=0$, $x_2=2$, $x_3=6$, etc. 
Each term is multiplying two numbers that are both
getting larger, so it is {\bf strictly increasing}. 
(f) This is similar to the previous one, but $x_0=x_1=0$, so it is just {\bf increasing}.
(g) This alternates between $-1$ and $1$, so it is {\bf non-monotonic}.  
(h) Each term subtracts from $1$ a smaller numbers than the last term, so it is {\bf strictly increasing}. 
(i) Each term adds to $1$ a smaller number than the last term, so it is {\bf strictly decreasing}.
\end{answer}
\end{exer}


% Totally unnecessary in our context.  But I'll leave it commented for now instead of
% just deleting it.  CAC, 5/21/14
%
%\begin{df}
%A sequence $\{x_n\} _{n=0} ^{+\infty}$ is said to be  {\em bounded}
%if eventually the absolute value of every term is less than or equal to a
%certain positive constant. A sequence that is not bounded is called {\em unbounded}.
%\index{bounded sequence}
%\index{unbounded sequence}
%\end{df}
%
%Proving that a sequence is unbounded involves showing that for any
%arbitrarily large positive real number, we can always find a term
%whose absolute value is greater than this real number.
%
%\begin{exa}
% Prove that the sequence $x_n = n!, n = 0, 1,
%2, \ldots$ is unbounded.
%\begin{pf}
%Let $M > 2$ be a large real number. Then $(\lceil M \rceil - 1)!>1$.
%Therefore $\{x_n\}$ is unbounded since
%$$x_{\lceil M\rceil}= \lceil M \rceil! = \lceil M \rceil(\lceil M \rceil - 1)! > \lceil M \rceil > M.$$
%\end{pf}
%\end{exa} 
%
%\begin{exa}
%Prove that the sequence $a_n = \dfrac{n + 1}{n},$ $n = 1, 2,
%\ldots,$ is bounded.
% \begin{pf}   
%Notice that if $n\geq 1$, then $1/n\leq 1$.
%Therefore 
%$$\displaystyle a_n=\frac{n+1}{n}=1+\frac{1}{n}\leq 2,$$ 
%so $a_n$ is bounded.
%\end{pf}
%\end{exa}

There are two types of sequences that come up often.  We will briefly discuss each.

\begin{df}\index{geometric progression}\index{geometric sequence}
A {\bf geometric progression} is a sequence of the form
$$ a, \ ar, ar^2,\ ar^3, \ ar^4, \ldots,$$
where $a$ (the {\bf initial term}) and $r$ (the {\bf common ratio}) are real numbers.
That is, a geometric progression is a sequence in which every term is
produced from the preceding one by multiplying it by a fixed number.
\end{df}

Notice that the first term can be written as $ar^0$, so like an array in many programming languages, the terms of a geometric progression 
are indexed starting at $0$. Thus, the $n$-th term is $ar^{n-1}$.
If $a=0$ then every term is $0$.
If $a r\neq 0$, we can find $r$ by dividing any term by the previous term.

\begin{exa}
Find the $11$-th term of the geometric progression
$$3, 6, 12, 24, \ldots . $$
 \begin{sol}  
Since this is a geometric progression, $a=3$ since the first term is always $a$. To determine $r$, we need to find the ratio between any two terms.
For instance, $24/12=2$ or $12/6=2$. So $r=2$ in this case.
Thus, the sequence is $\{3\cdot 2^k\}$, and the 11-th term is
$3\cdot 2^{10}=3*1024=3072$. 
\end{sol}
\end{exa}


\begin{exa}
Find the $35$-th term of the geometric progression
$$\frac{1}{\sqrt{2}},\  -2, \ \frac{8}{\sqrt{2}}, \ldots . $$
 \begin{sol}   
$a=\frac{1}{\sqrt{2}}$, and the common ratio is 
$r=-2/\frac{1}{\sqrt{2}} = -2\sqrt{2}$. 
Thus, the $n$-th term is $\frac{1}{\sqrt{2}}\left(-2\sqrt{2}\right)^{n-1}$.
Hence the $35$-th term is
$\frac{1}{\sqrt{2}}\left(-2\sqrt{2}\right)^{34} = \frac{2^{51}}{\sqrt{2}} =
1125899906842624\sqrt{2}.$
\end{sol}
\end{exa}

\begin{exer}
Find the $17$-th term of the geometric progression
$$-\frac{2}{3^{17}},\ \frac{2}{3^{16}},\ -\frac{2}{3^{15}}, \cdots . $$
\vspace*{1in}
\begin{answer}
You should have concluded that $a=-\frac{2}{3^{17}}$ and that
$r=\frac{2}{3^{16}}/\left(-\frac{2}{3^{17}}\right)=-3^{17}/3^{16}=-3$
(or you could have divided the second and third terms).
Then the $n$-th term is
$-\frac{2}{3^{17}}\left(-3\right)^{n-1}=\frac{2(-1)^{n}}{3^{18-n}}$
(Make sure you can do the algebra to get to this simplified form).
Finally, the $17$th term is
$\frac{2(-1)^{17}}{3^{18-17}}=-\dfrac{2}{3}$
\end{answer}
\end{exer}

\begin{exa}
The fourth term of a geometric progression is $24$ and its seventh
term is $192.$ Find its second term.
\begin{sol} 
We are given that $ar^3 = 24$ and $ar^6 = 192,$
for some $a$ and $r$. Clearly, $ar \neq 0,$ and so we find
$$\frac{ar^6}{ar^3} = r^3 = \frac{192}{24}  = 8.$$ 
Thus, $r = 2$.
Now, $a(2)^3 = 24,$ giving $a = 3.$ The second term is thus $ar =
6.$
\end{sol}
\end{exa}


\begin{exer}
The $6$-th term of a geometric progression is $20$ and the $10$-th
is $320$. Find  the absolute value of its third term.

\vspace{2in}
\begin{answer}
We are given that $ar^5 = 20$ and $ar^9 = 320$. 
Dividing, we can see that $r^4=16$.  Thus $r=\pm 2$.
(We don't have enough information to know which it is).
Since $ar^5=20$, we know that $a=20/r^5=\pm 20/32$.
So the third term is $ar^2=(\pm 20/32)(\pm 2)^2=\pm 80/32 = \pm 5/2$.
Thus $|ar^2|=5/2$.
\end{answer}
\end{exer}

\begin{df}\index{arithmetic progression}\index{arithmetic sequence}
An {\bf arithmetic progression} is a sequence of the form
$$ a, \ a+d, a+2d,\ a+3d, \ a+4d, \ldots,$$
where $a$ (the {\bf initial term}) and $d$ (the {\bf common difference}) are real numbers.
That is, an arithmetic progression is a sequence in which every term is
produced from the preceding one by adding a fixed number.
\end{df}

\begin{exa}
If $s_n=3n-7$, then $\{s_n\}$ is an arithmetic progression with $a=-7$ and $d=3$ (assuming we
begin with $s_0$).
\end{exa}

\begin{rem}
Notice that geometric progressions are essentially a discrete version of an exponential
function and arithmetic progressions are a discrete version of a linear function.
One consequence of this is that a sequence cannot be both of these unless it is the
sequence $a,a,a,\ldots$ for some $a$.
\end{rem}

\begin{exa}
Consider the sequence 
$$4, 7, 10, 13, 16, 19, 22,\ldots.$$  Assuming the pattern continues,
is this a geometric progression?  Is it an arithmetic progression?
\begin{sol}
It is easy to see that each term is 3 more than the previous term.  Thus, this is
an arithmetic progression with $a=4$ and $d=3$.   Clearly it is therefore not geometric. 
\end{sol}
\end{exa}


\begin{ques}
Tests like the SAT and ACT often have questions such as the following.
\begin{quote}
\begin{enumerate}
\item[23.]
Given the sequence of numbers 2, 9, 16, 23, what will the $8$th term of the 
sequence be?
(a) 60 (b) 58 (c) 49 (d) 51 (e) 56
\end{enumerate}
\end{quote}
\begin{lenum}
\item
What is the `correct' answer to this question?

\noindent
\answerlinetwo
\item 
Why did I put `correct' in quotes in the previous question?

\noindent
\answerlinetwo
\end{lenum}
\begin{answer}
\begin{lenum}
\item
The difference between the each of the first 4 terms of the sequence is 7, so it 
appears to be an arithmetic sequence.  Doing a little math, the correct answer
appears to be (d) 51.
\item
Although the sequence {\em appears to be} arithmetic, we cannot be certain that it is. 
If you are told it is arithmetic, then 51 is absolutely the correct answer.
Notice that the previous example specifically stated that you should assume that the
pattern continues.  This one did not.
Without being told this, the rest of the sequence could be anything. 
The $8$th term could be $0$ or $8,675,309$ for all we know.
Of the choices given, 51 is the most {\em obvious} choice, but any of the answers 
could be correct.  This is one reason I hate these sorts of questions on tests.

Although I think it is important to point out the flaw in these sorts of questions,
it is also important to conform to the expectations when answering such questions
on standardized tests.
In other words, instead of disputing the question (as some students might be inclined
to do), just go with the obvious interpretation.
\end{lenum}
\end{answer}
\end{ques}

Now let's see if you can correctly identify geometric and/or arithmetic 
sequences.

\newpage % better to have it on one page.
 
\begin{ques}
Determine whether or not the following sequences are geometric and/or arithmetic.
Explain your answer.
\begin{lenum}
\item
The sequence from Example~\ref{exer:5pow}.

\noindent 
\answerlinetwo
\item
The sequence from Example~\ref{exer:factrec}.

\noindent 
\answerlinetwo
\item
The sequence generated by \verb+ferzle(n)+ in Evaluate~\ref{eval:ferzle} 
on the non-negative inputs.

\noindent 
\answerlinetwo
\end{lenum}
\begin{answer}
(a) The closed form was $x_n=5^n$, which is clearly geometric (with $a=1$ and $r=5$) and not arithmetic.
(b) Since the solution for this one is $x_n=n!$, this is neither arithmetic or geometric.
(c) Since the sequence is essentially $f_n=2n+3$, with initial condition $f_0=3$, it is an
arithmetic sequence.  It is clearly not geometric.
\end{answer}
\end{ques}

\newpage

%%%%%%%%%%%%%%%%%%%%%%%%%%%%%%%%%%%%%%%%%%%%%%%%%%%%%%%%%%%%%%%%%%%
%%%%%%%%%%%%%%%%%%%%%%%%%%%%%%%%%%%%%%%%%%%%%%%%%%%%%%%%%%%%%%%%%%%
\section{Sums and Products}\label{section:sumProd}
When there is a need to add or multiply terms from a sequence, {\em summation notation} (or {\em sum notation})
and {\em product notation} come in handy.  We first introduce sum notation.

\begin{df}\index{a sum@$\sum$ (summation)}\index{sum notation}
Let $\{a_n\}$ be a sequence.  Then for $1\leq m\leq n$, where $m$ and $n$ are integers,
we define
$$\dsum _{k=m} ^n a_k=a_m+a_{m+1}+\cdots + a_n.$$
We call $k$ the {\bf index of summation} and $m$ and $n$ the {\bf limits} of the summation.
More specifically, $m$ is the {\bf lower limit} and $n$ is the {\bf upper limit}.
Each $a_k$ is a {\bf term} of the sum.
\end{df}

\begin{rem}
We often use $i$, $j$, and $k$ as index variables for sums, although
any letters can be used.
\end{rem}
\begin{exa}
We can express the sum $1 + 3 + 3^2 + 3^3 + \cdots + 3^{49}$ as
$$\sum_{i=0}^{49} 3^i.$$
(Recall that $3^0=1$, so the first term fits the pattern.)
\end{exa}

\begin{exer}
Write $1 + y + y^2 + y^3 + \cdots + y^{100}$ using sum notation.

\vspace*{.5in}
\begin{answer}
$\dsum_{i=0}^{100} y^i$
\end{answer}
\end{exer}

\begin{exa}
Write the following sum using sum notation.
$$1 - y + y^2 - y^3 + y^4 - y^5 +  \cdots - y^{99}+ y^{100}$$
\begin{sol}
This is a lot like the previous exercise, except that every other term is negative.
So how do we get those terms to be negative?  The standard trick relies on the fact that
$(-1)^i$ is $1$ if $i$ is even and $-1$ if $i$ is odd.  Thus, we can multiple each term by
$(-1)^i$ for an appropriate choice of $i$.  Since the odd powers are the negative ones, this
is easy:

$$\dsum_{i=0}^{100} (-1)^iy^i \ \ \ \ \ \mbox{ or } \ \ \ \ \  \dsum_{i=0}^{100} (-y)^i$$
\end{sol}
\end{exa}

\begin{rem}
You might be tempted to give the following solution to the previous problem:
$$\sum_{i=0}^{100} -y^i.$$
As we will see shortly, this is the same as
$$-\sum_{i=0}^{100} y^i,$$
which is not the correct answer.
The bottom line: Always use parentheses in the appropriate locations, especially when negative numbers are involved!
\end{rem}

\begin{exer}
Write $1 + y^2 + y^4 + y^6 + \cdots + y^{100}$ using sum notation.

\vspace*{.5in}
\begin{answer}
$\dsum_{i=0}^{50} (y^2)^i \ \ \ \ \ \mbox{ or } \ \ \ \ \ \dsum_{i=0}^{50} y^{2i}$
\end{answer}
\end{exer}

\begin{rem}
If you struggled understanding the two solutions to the previous example, it might be time to
review the basic algebra rules involving exponents.  We will just give a few of them here.
You can find more extensive lists in an algebra book or various reputable online sources.
We have already used the fact that if $x\neq 0$, then $x^0=1$.
In addition, if $x,a,b\in \reals$ with $x>0$, then 
$$
(x^a)^b=x^{ab}, \ \ \ \  
x^ax^b=x^{a+b}, \ \ \ \ 
(x^{-a})=\frac{1}{x^a}, \ \ \ \ \mbox{ and } \ \ \ \ 
x^{\frac{a}{b}}=\sqrt[b]{x^a}=\left(\sqrt[b]{x}\right)^a.$$
\end{rem}

As with sequences, we are often interested in obtaining {\em closed forms} for sums.
We will present several important formulas, along with a few techniques to find 
closed forms for sums.

\begin{exa}\label{exa:addOnes}
It should not be too difficult to see that 
$$\sum_{k=1}^{20} 1 = 20$$
since this sum is adding $20$ terms, each of which is $1$.
But notice that 
$$\sum_{k=0}^{19} 1 = \sum_{k=200}^{219} 1 = 20$$
since both of these sums are also adding $20$ terms, each of which is $1$.
In other words, if the variable of summation (the $k$) does not appear in the
sum, then the only thing that matters is how many terms the sum involves.
\end{exa}

\begin{exer}
Find each of the following.
\begin{lenum}
\item ${\displaystyle \sum_{k=5}^{6} 1 = \blankline{2} }$
\item ${\displaystyle \sum_{k=20}^{30} 1 = \blankline{2} }$
\item ${\displaystyle \sum_{k=1}^{100} 1 = \blankline{2} }$
\item ${\displaystyle \sum_{k=0}^{100} 1 = \blankline{2} }$
\end{lenum}
\begin{answer}
(a) $2$
(b)$11$
(c) $100$
(d) $101$
\end{answer}
\end{exer}

Hopefully you noticed that the previous example and exercise can be generalized as follows.

\begin{thm}\label{thm:sumones}
If $a,b\in\BBZ$, then 
$$\sum_{k=a}^b 1 = (b-a+1).$$
\begin{pf}
This sum has $b-a+1$ terms since there are that many number between $a$ and $b$, inclusive.
Since each of the terms is $1$, the sum is obviously $b-a+1$.
\end{pf}
\end{thm}

\begin{exa}
If we apply the previous theorem to the sums in Example~\ref{exa:addOnes},
we would obtain $20-1+1=20$, $19-0+1=20$, and $219-200+1=20$.
\end{exa}

Next is a simple theorem based on the distributive law that you learned in
grade school. 

\begin{thm}\label{thm:sumDist}
If $\{x_n\}$ is a sequence and $a$ is a real number, then 
$$\dsum _{k=m} ^n a\cdot x_k = a\dsum _{k=m} ^n x_k.$$
\end{thm}

\begin{exa}\label{exa:constantMult}
Using Theorems~\ref{thm:sumones} and~\ref{thm:sumDist}, we can see that
$$ \sum_{k=5}^{17} 4 = 4  \sum_{k=5}^{17} 1 = 4\cdot (17-5+1)=4\cdot 13 = 52.$$
\end{exa}

\begin{exer}
Find each of the following.
\begin{lenum}
\item ${\displaystyle \sum_{k=5}^{6} 5 = \blankline{4} }$
\item ${\displaystyle \sum_{k=20}^{30} 200 = \blankline{4} }$
\end{lenum}
\begin{answer}
(a) ${\displaystyle \sum_{k=5}^{6} 5 = 5 \sum_{k=5}^{6} 1= 5\cdot 2 = 10 }$.
(b) ${\displaystyle \sum_{k=20}^{30} 200 = 200\sum_{k=20}^{30} 1=200(30-20+1)=2200}$.
\end{answer}
\end{exer}

We can combine Theorems~\ref{thm:sumones} and~\ref{thm:sumDist} to obtain the following.
\begin{thm}\label{thm:sumconstant}
If $a,b\in\BBZ$ and $c\in\reals$, then 
$$\sum_{k=a}^b c = (b-a+1)c.$$
\begin{pf}
Using Theorem~\ref{thm:sumDist}, we have
$$\sum_{k=a}^b c = c\sum_{k=a}^b 1 = (b-a+1)c.$$
\end{pf}
\end{thm}

\begin{exa}
We can compute the sum from Example~\ref{exa:constantMult} by using Theorem~\ref{thm:sumconstant} to obtain
$$ \sum_{k=5}^{17} 4 = (17-5+1) 4 = 52.$$
Both ways of computing this sum are valid, so feel free to use whichever you prefer.
\end{exa}


\begin{exer}
Find each of the following.
\begin{lenum}
%\item ${\displaystyle \sum_{k=5}^{6} 5 = \blankline{2} }$
\item ${\displaystyle \sum_{k=20}^{30} 200 = \blankline{2} }$
\item ${\displaystyle \sum_{k=1}^{100} 9 = \blankline{2} }$
\item ${\displaystyle \sum_{k=0}^{100} 9 = \blankline{2} }$
\end{lenum}
\begin{answer}
Using Theorem~\ref{thm:sumconstant}, we get the following answers:
%(a) $2\cdot 5 = 10$.
(a) $(30-20+1)200=11*200=2200$.
(b) $900$
(c) $909$. Notice that this one has one more term than the previous one.  
The fact that the additional index is $0$ doesn't matter since it is adding $9$ for that term. 
\end{answer}
\end{exer}

\begin{eval}
Compute  ${\displaystyle \sum_{k=25}^{75} 10}$.
\begin{ssol}
This is just $10(75-25)=10*50=500$.
\end{ssol}
\evallinetwo
\begin{answer}
This solution contains an `off by one' error.  The correct answer is $10(75-25+1)=10*51=510$.
\end{answer}
\end{eval}

The following sum comes up often and should be committed to memory.
The proof involves a nice technique that adds the terms in the sum twice, 
in a different order, and then divides the result by two. This is known as Gauss' trick.
\begin{thm}\label{thm:sum-of-first-n-integ} 
If $n$ is a positive integer, then 
$$\sum_{k=1}^{n} k = \dfrac{n(n+1)}{2}.$$
\begin{pf} 
Let ${\displaystyle S=\sum_{k=1}^{n} k }$ for shorthand.  Then we can see
that
$$ S =  1 + 2 + 3 + \cdots + n $$
and by reordering the terms, 
$$ S =  n + (n - 1) +  \cdots + 1.$$
Adding these two quantities,
$$ \begin{array}{lcccccccc} S & = & 1 & + &  2 &  + &  \cdots & + & n \\
{S} & {=} & n & + & (n - 1) & + & \cdots & + & 1 \\
\hline
2S & =  & (n + 1) & + & (n + 1) & + & \cdots & + & (n + 1) \\
& = & n(n + 1), & & & & & 
\end{array}$$
since there are $n$ terms. 
Dividing by $2$, we obtain $S = \dis{\dfrac{n(n + 1)}{2}}$, as was to be proved.
\end{pf}
\end{thm}

\begin{exa}
$$\sum_{k=1}^{10}k=\frac{10(10+1)}{2}=\frac{10\cdot 11}{2}=55.$$
\end{exa}


\newpage % awkward break.

\begin{exer}
Compute each of the following.
\begin{lenum}
\item ${\displaystyle \sum_{k=1}^{20} k = \blankline{2} }$
\item ${\displaystyle \sum_{k=1}^{100} k = \blankline{2} }$
\item ${\displaystyle \sum_{k=1}^{1000} k = \blankline{2} }$
\end{lenum}
\begin{answer}
(a) $20\cdot 21/2=210$
(b) $100\cdot 101/2=5050$
(c) $1000\cdot 1001/2=500500$
\end{answer}
\end{exer}

\begin{eval}
Compute $\dsum_{k=1}^{30} k.$
\begin{ssolenum}
\item
${\displaystyle \sum_{k=1}^{30} k=29*30/2=435.}$

\noindent
\evallinetwo
\item
${\displaystyle \sum_{k=1}^{30} k=k\sum_{k=1}^{30}1 = k(30-1+1)=30k.}$

\noindent
\evallinetwo
\end{ssolenum}
\begin{answer}
\begin{solevalenum}
\item
Another example of the `off by one error'.  They are using the formula $n(n-1)/2$ 
instead of $n(n+1)/2$.
\item
This answer doesn't even make sense.  What is $k$ in the answer?  $k$ is just an
index of the summation.  The index should {\em never} appear in the answer.
The problem is that you can't pull the $k$ out of the sum since each term in the
sum depends on it.  
\end{solevalenum}
\end{answer}
\end{eval}

\begin{rem}
A common error is to think that the sum of the first $n$ integers is $n(n-1)/2$ 
instead of $n(n+1)/2$.  Whenever I use the formula, I double check my memory by
computing $1+2+3$.  In this case, $n=3$.  So is the correct answer $3\cdot 2/2=3$
or $3\cdot 4/2=6$?  Clearly it is the latter.  Then I know that the correct
formula is $n(n+1)/2$.  You can use any positive value of $n$ to check the formula.
I use $3$ out of habit.
\end{rem}

%\newpage % another awkward one

\begin{ques}
Is it true that $\displaystyle\sum_{k=0}^{n} k = \sum_{k=1}^{n} k = \dfrac{n(n+1)}{2}?$
Explain.

\noindent
\answerline
\begin{answer}
It is true.  The additional term that the sum adds is $0$, so the sum is the same whether
or not it starts at $0$ or $1$.
\end{answer}
\end{ques}


\begin{thm}
If $\{x_k\}$ and $\{y_k\}$ are sequences, then for any $n\in\BBZ^+$, 
$$\dsum_{i=1}^n x_i+y_i = \dsum_{i=1}^n x_i+\dsum_{i=1}^n y_i.$$
\begin{pf}
This follows from the commutative property of addition.
\end{pf}
\end{thm}

\begin{exa}
$$\dsum_{i=1}^{20} i+5 = \dsum_{i=1}^{20} i+ \dsum_{i=1}^{20} 5 = \frac{20\cdot 21}{2}+5\cdot 20 = 210+100=310.$$
\end{exa}

\begin{exer}
Compute the following sum

\noindent
$\dsum_{i=1}^{100} 2 - i = \blankline{5}.$
\begin{answer}
$\dsum_{i=1}^{100} 2 - i =\dsum_{i=1}^{100} 2 - \dsum_{i=1}^{100} i
=200 - \dfrac{100\cdot 101}{2} = 200-5050 = -4850.$
\end{answer}
\end{exer}


\begin{exer}
Prove that the sum of the first $n$ odd integers is $n^2$.

\vspace*{1.5in}
\begin{answer}
The sum of the first $n$ odd integers is

\noindent
$\dsum_{k=1}^{n}(2k-1)=
\dsum_{k=1}^{n}2k-\dsum_{k=1}^{n}1=
2\dsum_{k=1}^{n}k-\dsum_{k=1}^{n}1=
2\frac{n(n+1)}{2}-n=n^2+n-n=n^2.
$
\end{answer}
\end{exer}

The following example contains something called a {\em telescoping series}.
It demonstrates that evaluating a telescoping series is fairly simple.

\begin{exa}
Let $\{a_k\}$ be a sequence of real numbers.
Show that $\dsum_{i=1}^{n}(a_i-a_{i-1})=a_n-a_0$.

{\bf Proof: }
We can see that 
\begin{eqnarray*}
\dsum_{i=1}^{n}(a_i-a_{i-1})
&=&\left(\dsum_{i=1}^{n}a_i\right)-\left(\dsum_{i=1}^{n}a_{i-1}\right)\\
&=&(a_1+a_2+\cdots+a_{n-1}+a_n)-(a_0+a_1+a_2+\cdots +a_{n-1})\\
&=&a_1+a_2+\cdots+a_{n-1}+a_n-a_0-a_1-a_2-\cdots -a_{n-1}\\
&=&(a_1-a_1)+(a_2-a_2)+\cdots+(a_{n-1}-a_{n-1})+a_n -a_0\\
&=&a_n-a_0.\square
\end{eqnarray*}
\end{exa}



\begin{exa}
Given what we know so far, how can we compute the following:
$$ \sum_{k=50}^{100} k=?$$
It turns out that this is not that hard.  Notice that it is {\em almost} a sum we know.
We know how to compute ${\displaystyle \sum_{k=1}^{100} k}$, but that has too many terms.
Can we just subtract those terms to get the answer?  What terms don't we want?
Well, we don't want terms $1$ through $49$.  But that is just  ${\displaystyle \sum_{k=1}^{49} k}$.
In other words, 
\begin{eqnarray*}
\sum_{k=50}^{100} k&=&\sum_{k=1}^{100} k - \sum_{k=1}^{49} k\\
&=&\frac{100\cdot 101}{2} - \frac{49\cdot 50}{2}\\
&=&5050-1225=3825
\end{eqnarray*}
\end{exa}


\begin{exer}
Compute each of the following.
\begin{lenum}
\item ${\displaystyle \sum_{k=10}^{20} k = \blankline{5} }$
\item ${\displaystyle \sum_{k=21}^{40} k = \blankline{5} }$
\end{lenum}
\begin{answer}
(a) ${\displaystyle \sum_{k=10}^{20} k 
=\sum_{k=1}^{20} k - \sum_{k=1}^{9} k
=20\cdot 21/2 - 9\cdot 10/2 = 210 - 45 = 165}$.

(b) ${\displaystyle \sum_{k=21}^{40} k 
=\sum_{k=1}^{40} k - \sum_{k=1}^{20} k
=40\cdot 41/2 - 20\cdot 21/2 = 820-210 = 610}.$
\end{answer}
\end{exer}
\begin{eval}
Compute the following.
$$\sum_{k=30}^{100} k.$$
\begin{ssolenum}
\item
$${\displaystyle \sum_{k=30}^{100} k=
\sum_{k=1}^{100} k - \sum_{k=1}^{30} k=
100\cdot 101/2 - 30\cdot 31/2 = 5050 - 465 =4585
}$$

\noindent
\evallinetwo
\item
$${\displaystyle \sum_{k=30}^{100} k=
\sum_{k=1}^{100} k - \sum_{k=1}^{30} k=
99\cdot 100/2 - 29\cdot 30/2 = 4950 - 435 = 4515
}$$

\noindent
\evallinetwo
\item
$${\displaystyle \sum_{k=30}^{100} k=
\sum_{k=1}^{100} k - \sum_{k=1}^{29} k=
100\cdot 101/2 - 29\cdot 30/2 = 5050 - 435 =4615
}$$

\noindent
\evallinetwo
\end{ssolenum}
\begin{answer}
\begin{solevalenum}
\item
Another example of the off-by-one error.
The second sum should end at $29$, not $30$.
\item
This one has two errors, one of which is repeated twice.
It has the same error as the previous solution, but it also uses
the incorrect formula for each of the sums (the off-by-one error).
\item
This one is correct.
\end{solevalenum}
\end{answer}
\end{eval}

\begin{ques}
Explain why the following computation is incorrect.  Then explain why 
the answer is correct even with the error(s).
$$\sum_{k=30}^{100} k=
\sum_{k=1}^{100} k - \sum_{k=1}^{30} k=
100\cdot 101/2 - 29\cdot 30/2 = 5050 - 435 =4615
$$

\noindent
\answerlinefour
\begin{answer}
Two errors are made that cancel each other out.
The first error is that the second sum in the second step should go to $29$, not $30$.
But in the computation of that sum in the next step, 
the formula $n(n-1)/2$ is used instead of $n(n+1)/2$
(The correct formula was used for the first sum).
This is a rare case where an off-by-one error is followed by the opposite off-by-one error
that results in the correct answer.

It should be emphasized that even though the correct answer is obtained, 
{\em this is an incorrect solution}.
They obtained the correct answer by sheer luck.
\end{answer}
\end{ques}

%Perhaps the simplest cases are when we have a sum/product of the form
%$$(a_2-a_1)+(a_3-a_2)+\cdots + (a_n-a_{n-1})=a_n-a_1,  $$
%and
%$$ \dfrac{a_2}{a_1}\cdot  \dfrac{a_3}{a_2}\cdots  \dfrac{a_n}{a_{n-1}}= \dfrac{a_n}{a_1}, $$
%in which case we say that the sum or the product {\em telescopes}.
%The trick can be to recognize when a sum/product telescopes.
%
%In the following example we
%develop a formula for a geometric series by doing a little algebra so we can use the telescoping idea.

\newpage % Got a bad break at the top of this theorem for some reason.

\begin{thm}\label{thm:commonSums}
Let $n\in\BBZ^{+}$.  Then the following hold.
\begin{eqnarray*}
\dsum _{k=1}^n k^2 &=& \dfrac{n(n+1)(2n+1)}{6}\\
\dsum _{k=1}^n k^3 &=& \dfrac{n^2(n+1)^2}{4}\\
\dsum _{k=2}^n \dfrac{1}{(k-1)k}&=&\dfrac{1}{1\cdot 2} +
\dfrac{1}{2\cdot 3} + \dfrac{1}{3\cdot 4} + \cdots +
\dfrac{1}{(n-1)\cdot n} = \dfrac{n-1}{n}
\end{eqnarray*}
\end{thm}
We will prove Theorem~\ref{thm:commonSums} in the chapter on mathematical induction
since that is perhaps the easiest way to prove these results.
It is probably a good idea to attempt to commit the first two of these sums to memory since
they come up on occasion.

\begin{ques}
Why does the third formula from Theorem~\ref{thm:commonSums} have a lower index of $2$ 
(instead of $1$ or $0$, for instance)?

\noindent
\answerlinetwo
\begin{answer}
There are two ways to answer this.  The smart aleck answer is `because it is correct.'
But {\em why} is it correct with $2$, and couldn't it be slightly modified to work with $1$ or $0$?
The answer is {\em no} because if you plug $1$ or $0$ into $\frac{1}{(k-1)k}$, 
you get a division by $0$.  Hopefully I don't need to tell you that this is a bad thing.
\end{answer}
\end{ques}

\begin{exer}
Compute the following sum, simplifying as much as possible.

\noindent
$\hspace*{.25in}\dsum _{k=1}^n k^3+k=$
\vspace*{3.5in}
\begin{answer}
\begin{eqnarray*}
\dsum _{k=1}^n k^3+k &=& \dsum _{k=1}^n k^3+ \dsum _{k=1}^n k\\
&=&\frac{n^2(n+1)^2}{4}+\frac{n(n+1)}{2}\\
&=&\frac{n(n+1)}{2}\left(\frac{n(n+1)}{2}+1\right)\\
&=&\frac{n(n+1)}{2}\left(\frac{n^2+n+2}{2}\right)\\
&=&\frac{n(n+1)(n^2+n+2)}{4}
\end{eqnarray*}
\end{answer}
\end{exer}

Sometimes double sums are necessary to express a summation.
As a general rule, these should be evaluated from the inside out.

\begin{exa}
Evaluate the double sum 
$\dsum _{i=1}^n \dsum _{j=1}^n 1$.
\begin{sol}
We have
$\dsum _{i=1}^n \dsum _{j=1}^n 1 = \dsum _{i=1}^n n =n\cdot n = n^2.$
\end{sol}
\end{exa}

\begin{exer}
Evaluate the following double sums
\begin{lenum}
\item
$\dsum _{i=1}^n \ \dsum _{j=1}^i\ 1=$

\vspace*{1.6in}
\item
$\dsum _{i=1}^n \ \dsum _{j=1}^i\ j=$

\vspace*{1.5in}
\item
$\dsum _{i=1}^n \ \dsum _{j=1}^n\ ij=$

\vspace*{2in}
\end{lenum}
\begin{answer}
(a) 
$\dsum _{i=1}^n \dsum _{j=1}^i 1 = \dsum _{i=1}^n i = \dfrac{n(n+1)}{2}.$

\noindent
(b) $\dsum _{i=1}^n \ \dsum _{j=1}^i\ j
= \dsum _{i=1}^n \dfrac{i(i+1)}{2} = \cdots = \dfrac{n(n+1)(n+2)}{6}.$
(This one involves doing a little algebra, applying two formulas, and then doing a little more
algebra.  Make sure you work it out until you get this answer.)

\noindent
(c) 
$\dsum _{i=1}^n \ \dsum _{j=1}^n\ ij=
\dsum _{i=1}^n\left( i\dsum _{j=1}^n j\right)= 
\dsum _{i=1}^n\left( i\dfrac{n(n+1)}{2}\right)= 
\dfrac{n(n+1)}{2}\dsum _{i=1}^n i= 
\dfrac{n(n+1)}{2}\dfrac{n(n+1)}{2}= 
\dfrac{n^2(n+1)^2}{4}.$
\end{answer}
\end{exer}

There is a formula for the sum of a geometric sequence, sometimes referred to as a {\em geometric series}.
It is given in the next theorem.
\begin{thm}\label{thm:geom}
\index{geometric series} 
Let $x\neq 1$. Then
$$\dsum _{k=0}^n x^k =\dfrac{1-x^{n+1}}{1-x} \, \, \, \, \, \, \, \, \, \, 
\left(\mbox{ or } \, \, \, \, \, \dfrac{x^{n+1}-1}{x-1}\, \, \, \, \mbox{ if you prefer}\right).$$
\label{thm:sum-finite-geom-series}
\begin{pf}
First, let $S=\dsum _{k=0}^n x^k$.
Then $$xS = x\dsum _{k=0}^n  x^k = \dsum _{k=0}^n x^{k+1}= \dsum _{k=1}^{n+1}  x^{k}.$$
So
\begin{eqnarray*}
xS - S &=&\dsum _{k=1}^{n+1} x^{k} - \dsum _{k=0}^n x^k\\
&=&(x_1+x_2+\ldots+x_n+x_{n+1}) - (x_0+x_1+\ldots+x_n)\\
&=&x^{n+1}-x^0=x^{n+1}-1.
\end{eqnarray*}
So we have $(x-1)S=x^{n+1}-1$, so $S=\frac{x^{n+1}-1}{x-1}$, since $x\neq 1$.
\end{pf}
\end{thm}

\begin{exa}
$$\dsum _{k=0}^n 3^k =\dfrac{1-3^{n+1}}{1-3}=\dfrac{1-3^{n+1}}{-2}=\dfrac{3^{n+1}-1}{2}.$$
\end{exa}

\begin{exa}
$$\dsum _{k=0}^n \dfrac{1}{5^k} 
=\dsum _{k=0}^n \dfrac{1^k}{5^k} 
=\dsum _{k=0}^n \left(\dfrac{1}{5}\right)^k
=\dfrac{1-(1/5)^{n+1}}{1-1/5}
=\dfrac{1-1/(5^{n+1})}{4/5}
=\dfrac{5}{4}\left(1-\dfrac{1}{5^{n+1}}\right)
=\dfrac{5}{4}-\dfrac{1}{4\cdot5^{n}}.$$
\end{exa}

\begin{exer}
Find the sum of the following geometric series.

\noindent
$1 + 3 + 3^2 + 3^3 + \cdots + 3^{49}=$ 
\vspace*{1in}
\begin{answer}
$\frac{3^{50} - 1}{2}=358948993845926294385124$.
\end{answer}
\end{exer}

\begin{exer}
Find the sum of the following geometric series.

\noindent
$1-2+4-8+\cdots-2^{33}+2^{34}=$

\vspace*{1in}
\begin{answer}
This is equivalent to $\sum_{k=0}^{34}(-2)^k$,
 so the summation is 
 $(1-(-2)^{35})/(1-(-2))=(1-(-1)^{35}2^{35})/3=(1+2^{35})/3=11453246123$.
\end{answer}
\end{exer}

\begin{exer}
Find the sum of the following geometric series.
Assume $y\neq 1$.
\begin{lenum}
\item 
$1 + y + y^2 + y^3 + \cdots + y^{100}=$

\vspace*{1in}
\item 
$1 - y + y^2 - y^3 + y^4 - y^5 +  \cdots - y^{99}+ y^{100}=$

\vspace*{1.in}
\item 
$1 + y^2 + y^4 + y^6 + \cdots + y^{100}=$

\vspace*{1.1in}
\end{lenum}
\begin{answer}
 (a) $\frac{1 - y^{101}}{1 - y}$ or $\frac{y^{101}-1}{y-1}$ (We won't give the alternatives for the rest.
 If your answer differs, do some algebra to make sure it is equivalent.)
 (b) $\frac{1 -(-y)^{101}}{1 - (-y)} = \frac{1 + y^{101}}{1 + y}$
 (c) $\frac{1 - y^{102}}{1 - y^2}$.
\end{answer}
\end{exer}

\begin{cor}
Let $N\geq 2$ be an integer.  Then 
$$x^N-1 = (x-1)(x^{N-1}+x^{N-2}+\cdots + x +1).$$
\begin{pf}
Plugging $N=n+1$ in the formula from Theorem~\ref{thm:geom} and doing a little algebra
yields the formula.
\end{pf}
\end{cor}

\newpage % bad break on this one? Weird. This fixes it.
\begin{exa}
We can see that
\begin{eqnarray*}
x^2-1&=&(x-1)(x+1)\\
x^3-1&=&(x-1)(x^2+x+1), \mbox{ and}\\
x^4-1&=&(x-1)(x^3+x^2+x+1).
\end{eqnarray*}
\end{exa}

\begin{exer}
Factor $x^5-1$.

\noindent
$x^5-1=$\blankline{3}
\begin{answer}
$x^5-1=(x-1)(x^4+x^3+x^2+x+1)$.
\end{answer}
\end{exer}

%\begin{rem}
%More important than remembering the {\em formula} above is remembering the {\em method} of how this formula was obtained.  As you work with sums more, you will start to see some of the tricks that come in handy often.
%There may be ``one ring to rule them all'', but there is not one technique that always works when finding
%closed forms for sums.
%\end{rem}

Let's use the technique from the proof of Theorem~\ref{thm:geom} in the special case where $x=2$.

\begin{fib}
Find the sum $$2^0+2^1 + 2^2 + 2^3 + 2^4 + \cdots + 2^n.$$
 \begin{sol}   
We could just use the formula from Theorem~\ref{thm:geom}, but that would be boring.
Instead, let's work it out. 
Let $S = 2^0+2^1 + 2^2 + 2^3 + \cdots + 2^{n}.$
Then $2S=\blankline{2}$.  Notice $S$ and $2S$ have most of the same terms,
except $S$ has \blankline{1} that $2S$ doesn't have and 
$2S$ has \blankline{1} that $S$ doesn't have.
Therefore,
$$
\begin{array}{lllllllllllllll}
S = 2S - S &= &   &   & (2^1 & + & 2^2 & + & 2^3 & + & \cdots & + & 2^{n}  & + & 2^{n+1})\\
   &  &-( 2^0 & + & 2^1 & + & 2^2 & + & 2^3 & + & \cdots & + & 2^{n} ) &   &  \\
&=&\multicolumn{8}{l}{\blankline{2}}\\
&=&\multicolumn{8}{l}{2^{n+1}-1.}\\
\end{array}
$$
Thus, 
$\displaystyle\sum_{k=0}^n 2^k = 2^{n+1}-1.$
\end{sol}
\begin{answer}
$2^1+2^2 + 2^3  + \cdots + 2^{n+1}$;
$2^0$; 
$2^{n+1}$;
$2^{n+1}-2^0$
\end{answer}
\end{fib}

Since powers of $2$ are very prominent in computer science,
you should definitely commit the formula from the previous example to memory.

Together, Theorems~\ref{thm:sumDist} and~\ref{thm:sum-finite-geom-series} imply the following:

\begin{thm}\label{thm:geoser2}
Let $r\neq 1$.  Then 
$\displaystyle\sum_{k=0}^{n}ar^k  =\frac{a-ar^{n+1}}{1-r}.$
\end{thm}
\newpage

\begin{fib}
Use  Theorems~\ref{thm:sumDist} and~\ref{thm:sum-finite-geom-series} to prove Theorem~\ref{thm:geoser2}.
\begin{pf}
It is easy to see that
\begin{eqnarray*}
\dsum_{k=0}^{n}ar^k &=&\\[.5in]
&=&\\[.5in]
&=&\frac{a-ar^{n+1}}{1-r}.
\end{eqnarray*}
\end{pf}
\begin{answer}
$a\dsum_{k=0}^{n}r^k$;
$a\frac{1-r^{k+1}}{1-r}$.
\end{answer}
\end{fib}

\begin{exer}
Prove Theorem~\ref{thm:geoser2} {\em without using  
Theorems~\ref{thm:sumDist} and~\ref{thm:sum-finite-geom-series}.}
In other words, mimic the proof of Theorem~\ref{thm:sum-finite-geom-series}.

\vspace*{3.5in}
\begin{answer}
Let $S = a + ar + ar^2 + \cdots + ar^{n}.$ 
Then $rS = ar + ar^2 + \cdots + ar^{n+1}Z$, so 
\begin{eqnarray*}
S - rS &=& a + ar + ar^2 + \cdots + ar^{n} - ar - ar^2 - \cdots - ar^{n+1}\\
&=& a - ar^{n+1}.
\end{eqnarray*}
From this we deduce that $$S = \frac{a - ar^{n+1}}{1 - r},$$
that is,
$$ \sum_{k=0}^{n}ar^k= \frac{a - ar^{n+1}}{1 - r}.$$
\end{answer}
\end{exer}


Notice that if $|r| < 1$ then $r^n$ gets closer to $0$ the larger $n$ gets.
More formally, if $|r|<1$,  ${\displaystyle \lim_{n\rightarrow \infty} r^n=0}$.
This implies the following (which we will not formally prove beyond what
we have already said here).

\begin{thm}
Let $|r| < 1$.  Then 
$$\sum_{k=0}^{\infty} ar^k  =  \frac{a}{1 - r}.$$
\end{thm}

\begin{exa} 
A fly starts at the origin and goes $1$ unit up, $1/2$ unit right, $1/4$ unit down, $1/8$ unit left,
$1/16$ unit up, etc., {\em ad infinitum.} In what coordinates does
it end up?
\begin{sol}
Its $x$ coordinate is
$$
\frac{1}{2} - \frac{1}{8} + \frac{1}{32} - \cdots=
\frac{1}{2}\left(-\frac{1}{4}\right)^0 + \frac{1}{2}\left(-\frac{1}{4}\right)^1 + \frac{1}{2}\left(-\frac{1}{4}\right)^2 + \cdots
= \frac{\frac{1}{2}}{1 - \frac{-1}{4}} = \frac{2}{5}.$$ 
Its $y$ coordinate is
$$1 - \frac{1}{4} + \frac{1}{16} - \cdots =
\left(-\frac{1}{4}\right)^0 +\left(-\frac{1}{4}\right)^1 + \left(-\frac{1}{4}\right)^2 +\cdots =
 \frac{1}{1 - \frac{-1}{4}} = \frac{4}{5}.$$Therefore, the
fly ends up in $\left(\frac{2}{5}, \frac{4}{5}\right).$
\end{sol}
\end{exa}


The following infinite sums are sometimes useful.
\begin{thm}
Let $x\in \BBR$.  The following expansions hold:
$$
\begin{array}{lllll} 
\sin x &=&\dsum _{n=0} ^{\infty}\dfrac{(-1)^nx^{2n+1}}{(2n+1)!}&=& x - \dfrac{x^3}{3!} + \dfrac{x^5}{5!} - \cdots + (-1)^n\dfrac{x^{2n + 1}}{(2n + 1)!} +  \cdots\\[3ex]
\cos x &=&\dsum _{n=0} ^{\infty}\dfrac{(-1)^nx^{2n}}{(2n)!} &=&1 -
\dfrac{x^2}{2!} + \dfrac{x^4}{4!} - \cdots +(-1)^n\dfrac{x^{2n}}{(2n)!} +  \cdots \\[3ex]
e^x&=&\dsum _{n=0} ^{\infty}\dfrac{x^n}{n!}
&=& 1 + x + \dfrac{x^2}{2!} +  \dfrac{x^3}{3!} + \cdots + \dfrac{x^{n}}{n!}+  \cdots\\[3ex]
\dfrac{1}{1-x}&=&\dsum _{n=0} ^{\infty}x^n&=&1+x+x^2+x^3+\cdots, \mbox{ if } |x|<1
\end{array}
$$
\label{thm:maclaurin}
\end{thm}


{\em Product notation} is very similar to sum notation, except we multiply the terms instead of 
adding them.

\begin{df}
\index{a sumb@$\prod$ (product)}
Let $\{a_n\}$ be a sequence.  Then for $1\leq m\leq n$,
where $m$ and $n$ are integers,
we define
$$\prod _{k=m} ^n a_k= a_{m}a_{m+1}\cdots a_n.$$
As with sums, we call $k$ the {\bf index} and 
$m$ and $n$ the {\bf lower limit} and {\bf upper limit}, respectively.
\end{df}

\begin{exa}
Notice that $\displaystyle n! = \prod_{k=1}^n k$.
\end{exa}

\begin{rem}
An alternative way to express the variable and limits of sums and products is
$$ \dsum_{m\leq k \leq n} a_k \qquad \mbox{ instead of } \qquad \dsum_{k=m}^{n}a_k$$
and
$$\qquad \prod_{m\leq k \leq n}a_k \qquad\mbox{ instead of } \qquad \prod_{k=m}^{n} a_k$$
\end{rem}

%%--------------------------------------------------------------------------------------------
%%--------------------------------------------------------------------------------------------

\newpage

\section{Matrices and Matrix Operations}\label{sec:matrix}
\subsection{Definitions}\label{subsec:matDef}
\begin{df}\index{matrix}
An  $m\times n$ (read $m$ by $n$) {\bf matrix} $A$ with $m$ rows and $n$ columns with entries
over $\BBR$ is a rectangular array of the form
$$A = \begin{bmatrix}a_{11} & a_{12} & \cdots & a_{1n} \cr a_{21} & a_{22} & \cdots & a_{2n} \cr
\vdots & \vdots & \cdots & \vdots \cr a_{m1} & a_{m2} & \cdots &
a_{mn} \cr \end{bmatrix},$$ where $\forall (i, j) \in \{1, 2, \ldots
, m \}\times \{1, 2, \ldots , n\}, \ \ a_{ij}\in \BBR$.

\noindent
The number of rows and columns of a matrix is referred to
as its {\bf dimensions}.
\end{df}

\begin{exa}
$A = \begin{bmatrix}0 & 4 & 1 \cr 1 & 2 & 3 \cr\end{bmatrix}$ 
is a $2\times 3$ mtrix and 
$B = \begin{bmatrix}-2 & 1 \cr 1 & 2 \cr 0 & 3 \cr\end{bmatrix}$ 
is a $3\times 2$ matrix.
\end{exa}
\begin{rem}
As a shortcut, we use the notation $A = [a_{ij}]$ to denote
the matrix $A$ with entries $a_{ij}$. 
\end{rem}
\begin{rem}
In most programming languages, matrices are
indexed starting with $0$ instead of $1$. So the row indices 
go from $0$ to $m-1$ and the column indices go from $0$ to $n-1$.
There is a very good technical reason for this, but we will not worry
about that for now.
We will stick with a starting index of $1$ in these notes
since in mathematics that is more common.
\end{rem}

\begin{df}
To refer to a specific entry of a matrix, we use the
notation $a_{ij}$ (without square brackets).
Sometimes we use a comma, as in $a_{i,j}$,
to separate the two subscripts.
Alternatively, the notation $A[i,j]$ or
$A[i][j]$ is used, 
especially in programming languages.
\end{df}
\begin{exa}
If $\displaystyle A = \begin{bmatrix}0 & -1 & 1 \cr 1 & 2 & 3 \cr\end{bmatrix}$, then $a_{2,3}=A[2,3]=A[2][3]=3$.
Also, $A[1,2]=-1$.
\end{exa}
\begin{exa}
Write out explicitly the $4\times 4$ matrix $A = [a_{ij}]$ where
$a_{ij} = i^2 - j^2$.
\begin{sol} 
$A = \begin{bmatrix} 1^2 - 1^1 & 1^2 - 2^2 & 1^2 -
3^2 & 1^2 - 4^2 \cr 2^2 - 1^2 & 2^2 - 2^2 & 2^2 - 3^2 & 2^2 - 4^2
\cr 3^2 - 1^2 & 3^2 - 2^2 & 3^2 - 3^2 & 3^2 - 4^2 \cr 4^2 - 1^2 &
4^2 - 2^2 & 4^2 - 3^2 & 4^2 - 4^2 \cr
\end{bmatrix} =
\begin{bmatrix}   0 & -3& -8 & -15 \cr
 3 & 0 & -5 & -12 \cr
 8 & 5 & 0 & -7 \cr
 15 & 12 & 7 & 0 \end{bmatrix}. $
 \end{sol}
\end{exa}

\begin{exer}
Write out explicitly the $3\times 3$ matrix $A = [a_{ij}]$ where
$a_{ij} = i^j$.
\vspace*{1.5in}
\begin{answer} $A = \begin{bmatrix} 1 & 1 & 1 \cr
2 & 4 & 8 \cr 3 & 9 & 27 \cr \end{bmatrix}$  
\end{answer}
\end{exer}

\begin{df}
We denote by
$\mat{m\times n}{\BBR}$ the set of all $m\times n$ matrices with
entries over $\BBR$. 
Since we will always be working over $\BBR$ in this book, 
we will shorthand this notation to $\mats{m\times n}$.
$\mats{n\times n}$ is, in particular,
the set of all square matrices of size $n$.
\end{df}

\begin{df}
The $m\times n$ {\bf zero matrix}, ${\bf 0}_{m\times
n}\in\mats{m\times n}$, is the matrix with $0$'s
everywhere,
$${\bf 0}_{m\times n} = \begin{bmatrix}0 & 0 & 0 & \cdots & 0 \cr 0 & 0 & 0 & \cdots & 0 \cr 0 & 0 & 0
 & \cdots & 0 \cr
\vdots & \vdots & \vdots & \cdots & \vdots \cr 0 & 0
& 0 & \cdots & 0 \cr\end{bmatrix}.$$ When $m = n$ we
write ${\bf 0}_n$ as a shortcut for ${\bf 0}_{n\times
n}$.\index{matrix!zero}
\end{df}
\begin{df} \index{matrix!identity}
The $n\times n$ {\bf identity matrix}, 
${\bf I}_n\in\mats{n\times n}$, is the matrix with $1$'s on the  main diagonal and
$0$'s everywhere else,
$${\bf I}_n = \begin{bmatrix}1 & 0 & 0 & \cdots & 0 \cr 0 & 1 & 0 & \cdots & 0 \cr 0 & 0 & 1
& \cdots & 0 \cr \vdots & \vdots & \vdots & \cdots & \vdots
\cr 0 & 0 & 0 & \cdots & 1
\cr\end{bmatrix}.$$
\end{df}

\begin{exer}
Explicitly write out the entries of the matrices 
${\bf 0}_{3\times 4}$ and ${\bf I}_5$.
\vspace*{2in}
\begin{answer}
${\bf 0}_{3\times 4}=
\begin{bmatrix}0 & 0 & 0 & 0 \cr 
0 & 0 & 0 & 0 \cr 
0 & 0 & 0 & 0 \cr \end{bmatrix}$
and 
 ${\bf I}_5=
\begin{bmatrix}
1 & 0 & 0 & 0 & 0\cr 
0 & 1 & 0 & 0 & 0\cr 
0 & 0 & 1 & 0 & 0\cr 
0 & 0 & 0 & 1 & 0\cr 
0 & 0 & 0 & 0 & 1\cr\end{bmatrix}$
\end{answer}
\end{exer}

\begin{df}[Matrix Addition and Multiplication of a Matrix by a Scalar]\index{matrix!addition}\index{matrix!scalar multiplication} \ \\
Let $A = [a_{ij}], B = [b_{ij}] \in
\mats{m\times n}$ and $\alpha\in \BBR$. 
\begin{itemize}
\item
The matrix $\alpha A$ is the matrix 
$C \in \mats{m\times n}$ with entries
$C = [c_{ij}]$ where $c_{ij} = \alpha a_{ij}$.
\item
The matrix $A + B$ is the matrix $C \in \mats{m\times n}$ with entries
$C = [c_{ij}]$ where $c_{ij} = a_{ij}+b_{ij}$.
\end{itemize}
\end{df}
\begin{exa}
Let $\displaystyle A = \begin{bmatrix}2 & -3 & 1 \cr 4 & -7 & 3 \cr\end{bmatrix}$ and 
$\displaystyle B = \begin{bmatrix}0 & 3 & -5 \cr 6 & 0 & 4 \cr\end{bmatrix}$.
Then 
 $$ 3 A = \begin{bmatrix}3\times2 & 3\times-3 & 3\times1 \cr 3\times4 & 3\times-7 & 3\times3 \cr\end{bmatrix}
=\begin{bmatrix}6 & -9 & 3 \cr 12 & -21 & 9 \cr\end{bmatrix}$$
and
$$A+B=\begin{bmatrix}2+0 & -3+3 & 1-5 \cr 4+6 & -7+0 & 3+4 \cr\end{bmatrix}
=\begin{bmatrix}2 & 0 & -4 \cr 10 & -7 & 7 \cr\end{bmatrix}.$$
\end{exa}

\begin{exer}
Let $A = \begin{bmatrix} 1 & 1 \cr -1 & 1 \cr 0 & 2 \cr
\end{bmatrix}$ and $B = \begin{bmatrix} -1 & 1 \cr 2 & 1 \cr 0 & -1 \cr
\end{bmatrix}$.
Compute $A+2B$.

\vspace*{1.5in}
\begin{answer}
$\displaystyle A + 2B = \begin{bmatrix} -1 & 3 \cr 3 & 3 \cr 0 & 0
\end{bmatrix}$
\end{answer}
\end{exer}


\begin{exer}
Let $a,b,c\in\BBR$ and let 
$$M = \begin{bmatrix}a & -2a & c \cr  0 & -a & b \cr a + b & 0 &
-1 \cr\end{bmatrix},\ \ \ \ N = \begin{bmatrix}1 & 2a & c \cr  a & b
- a & -b \cr a - b & 0 & -1 \cr \end{bmatrix}.$$
Compute $M+N$ and $2M$.
\vspace*{2.5in}
\begin{answer}
$M + N = \begin{bmatrix}a + 1 & 0 & 2c \cr  a &  b - 2a & 0 \cr 2a
& 0 & -2 \cr\end{bmatrix}, \ \ \ 2M = \begin{bmatrix}2a & -4a & 2c
\cr 0 & -2a & 2b \cr 2a + 2b & 0 & -2 \cr\end{bmatrix}$
\end{answer}
\end{exer}


The following properties of matrices can easily be proved by 
applying the definitions of matrix addition and scalar multiplication.

\begin{thm}Let $A, B, C\in \mats{m\times n}$ and $\alpha, \beta \in \BBR$. Then
\begin{enumerate}
\item $\mats{m\times n}$ is closed under matrix addition
and scalar multiplication
\begin{equation}\label{m:closure} A + B \in \mats{m\times n}, \
\ \ \ \alpha A \in \mats{m\times n}
\end{equation}
\item Addition of matrices is commutative
\begin{equation}\label{m:commutative} A + B = B + A
\end{equation}
\item Addition of matrices is associative
\begin{equation}\label{m:associative} A + (B + C) = (A+B)+C
\end{equation}
\item  The zero matrix is the additive identity
\begin{equation}\label{m:existence0}A+ {\bf 0}_{m\times n}=A
\end{equation}
\item Additive inverse
\begin{equation}\label{m:existence_additive_inverse} 
A + (-A) = (-A) + A = {\bf 0}_{m\times n}
\end{equation}
\item Distributive law
\begin{equation}\label{m:distributive_law1} \alpha (A + B) =
\alpha A + \alpha B
\end{equation}
\item Distributive law
\begin{equation}\label{m:distributive_law2} (\alpha + \beta)A =
\alpha A + \beta A
\end{equation}
\item Scalar multiplicative identity
\begin{equation}\label{m:1A} 1A = A
\end{equation}
\item   
\begin{equation}\label{m:abA} \alpha (\beta A) =
(\alpha\beta)A
\end{equation}

\end{enumerate}
\end{thm}


\begin{exer}
Determine $x$ and $y$ such that $$\begin{bmatrix} 3 & x & 1 \cr 1
& 2 & 0 \cr
\end{bmatrix} + 2\begin{bmatrix} 2 & 1 & 3 \cr 5 & x & 4 \cr\end{bmatrix} = \begin{bmatrix}
7 & 3 & 7 \cr 11 & y & 8 \cr\end{bmatrix}.$$
\vspace*{1.5in}
\begin{answer} $x = 1$ and $y = 4$.  \end{answer}
\end{exer}

\subsection{Matrix Multiplication}\label{subsec:matMul}

Matrix multiplication is more straightforward that the following definition
makes it seem. Once you have a little practice with it, it comes very naturally.
\begin{df}\index{matrix!multiplication}
Let $A = [a_{ij}] \in \mats{m\times n}$ and $B = [b_{ij}] \in
\mats{n\times p}$. Then the matrix product $AB$ is defined as
the matrix $C = [c_{ij}]\in \mats{m\times p}$ with entries
$c_{ij} = \sum _{l=1} ^{n} a_{il} b_{lj}$:
$$\begin{bmatrix}a_{11} & a_{12} & \cdots & a_{1n} \cr a_{21} & a_{22} & \cdots & a_{2n}
\cr \vdots & \vdots & \cdots & \vdots  \cr  {\red a_{i1}} & {\red
a_{i2}} & {\red \cdots} & {\red a_{in}}\cr  \vdots & \vdots &
\cdots & \vdots \cr a_{m1} & a_{m2} & \cdots & a_{mn} \cr
\end{bmatrix}\begin{bmatrix}b_{11} &\cdots & {\blue b_{1j}} & \cdots & b_{1p}
\cr b_{21} &\cdots & {\blue b_{2j}} & \cdots & b_{2p} \cr \vdots
&\cdots & {\blue \vdots} & \cdots & \vdots \cr b_{n1} &\cdots &
{\blue b_{nj}} & \cdots & b_{np} \cr
\end{bmatrix} = \begin{bmatrix} c_{11} & \cdots & c_{1p} \cr c_{21} & \cdots & c_{2p} \cr
 \vdots & \cdots & \vdots \cr
 \cdots & {\green c_{ij}} & \cdots \cr \vdots & \cdots & \vdots \cr c_{m1} & \cdots & c_{mp} \cr\end{bmatrix}.$$
\end{df}
\begin{rem}
\begin{enumerate}
\item
Observe that we use juxtaposition rather than a special symbol to
denote matrix multiplication. This will simplify notation.
\item
In order
to obtain the $ij$-th entry of the matrix $AB$ we multiply
element-wise the $i$-th row of $A$ by the $j$-th column of $B$.
\item
Observe that $AB$ is a $m\times p$ matrix, and that in order
to multiply two matrices, the number of columns of the first 
matrix must be equal to the number of rows of the second one.
\end{enumerate}
\end{rem}


\begin{exa}
Let $M = \begin{bmatrix}1 & 2  \cr 3 & 4\cr\end{bmatrix}$ and $N =
\begin{bmatrix}5 & 6  \cr 7& 8 \cr \end{bmatrix}$. Then
$$MN =\begin{bmatrix}1 & 2  \cr 3 & 4\cr\end{bmatrix}\begin{bmatrix}5 & 6  \cr 7& 8 \cr \end{bmatrix}
= \begin{bmatrix} 1\cdot 5 + 2\cdot 7 &  1\cdot 6 + 2\cdot 8 \cr
3\cdot 5 + 4\cdot 7 & 3\cdot 6 + 4 \cdot 8 \cr\end{bmatrix} =
\begin{bmatrix}19 & 22  \cr 43 & 50 \cr \end{bmatrix},$$ and
$$NM =
\begin{bmatrix}5 & 6  \cr 7& 8 \cr
\end{bmatrix}\begin{bmatrix}1 & 2  \cr 3 & 4\cr\end{bmatrix}=
\begin{bmatrix}5\cdot 1 + 6\cdot 3 &  5\cdot 2 + 6\cdot 4 \cr 7\cdot 1 + 8\cdot 3 & 7\cdot 2 + 8\cdot 4\cr
\end{bmatrix} = \begin{bmatrix}23 & 34  \cr 31& 46 \cr \end{bmatrix}.$$
Notice that matrix multiplication is not necessarily commutative!
\end{exa}

\begin{exer}
Consider the matrix
$$A = \begin{bmatrix} 2 & 1 & 3 \cr
 0 & 1 & 1 \cr
  4 & 4 & 0 \cr
  \end{bmatrix}.$$ 
Compute $AA$.
\vspace*{1.25in}
\begin{answer}
 $\begin{array}{llll}AA & = &
\begin{bmatrix} 2 & 1 & 3 \cr
 0 & 1 & 1 \cr
  4 & 4 & 0 \cr
  \end{bmatrix}\begin{bmatrix} 2 & 1 & 3 \cr
 0 & 1 & 1 \cr
  4 & 4 & 0 \cr
  \end{bmatrix} \vspace{2mm}
   =  \begin{bmatrix} 16 & 15 & 7 \cr
 4 & 5 & 1 \cr
  8 & 8 & 16 \cr
  \end{bmatrix}
 \end{array} $
\end{answer}
\end{exer}

\begin{exa}
Notice that $$\begin{bmatrix} {1} & {1} & {1} \cr
 {1}  & {1} & 1 \cr
{1}  & {1} & {1} \cr
\end{bmatrix}
\begin{bmatrix}
{2}  & -{1} & -{1} \\ \vspace{1mm} -{1} & {2} & -{1}
\\ \vspace{1mm}
-{1}  & -{1} & {2} \\
\end{bmatrix}  = \begin{bmatrix}
0  & 0 & 0 \cr \vspace{1mm} 0  & 0 & 0 \cr \vspace{1mm} 0  & 0 & 0
\cr
\end{bmatrix}.$$
Observe then that  the product of two non-zero matrices may be
the zero matrix.
\end{exa}

\newpage

\begin{exer}
Let $a,b,c\in \BBR$,
 $A = \begin{bmatrix}1 & 0 & 0 \cr 1 & 1 & 0 \cr 1 & 1 & 1 \cr
\end{bmatrix},  \mbox{ and } B = \begin{bmatrix}
a & b & c \cr c & a & b \cr b & c & a \cr \end{bmatrix}$. 
Find
$AB$ and $BA$.
\vspace*{3.25in}
\begin{answer}  
$AB = \begin{bmatrix} a & b & c \cr c+a & a + b & b + c
\cr a + b + c & a + b + c & a + b + c \cr
 \end{bmatrix}, \ BA = \begin{bmatrix} a + b + c & b + c & c
 \cr a + b + c & a + b & b \cr a + b + c & c + a & a \cr
\end{bmatrix}
 $ 
 \end{answer}
\end{exer}


Even though matrix multiplication is not necessarily commutative,
it is associative.
\begin{thm}
If $A \in 
\mats{m\times n },
B \in \mats{n\times r}, \mbox{ and }
C \in \mats{r \times s}$ then $$(AB)C = A(BC),$$
i.e., matrix multiplication is associative.
\begin{pf} 
To prove this we only need to consider the $ij$-th entry of each side,
appeal to the associativity of mulitplication of real numbers, and
verify that both sides are indeed equal to
$$\sum _{k = 1} ^{n}\sum _{k' = 1} ^{r} a_{ik}b_{kk'}c_{k'j}.$$
\end{pf}
\end{thm}

\begin{df}
Let $A \in \mats{n\times n}$.
The notation $A^k$ has the obvious
meaning of $k$ copies of $A$ multiplied together.
We can define it more formally as $A^k=AA^{k-1}$ if $k>1$.
\end{df}


\begin{rem}
By virtue of associativity, a square matrix commutes with its
powers, that is, if $A\in\mats{n\times n}$, and 
$r,s\in\BBN$, then $(A^r)(A^s) = (A^s)(A^r) = A^{r + s}$.
\end{rem}

\begin{exer}
Let $M = \begin{bmatrix} 1 & -1  \cr -1 & 1 \cr\end{bmatrix}$.
Find $M^6$.
(Hint: If you are clever, you can do this with 3 multiplications
instead of 5!)
\vspace*{2.5in}
\begin{answer} $ \begin{bmatrix} 32 & -32  \cr -32 & 32 \cr\end{bmatrix}$.\end{answer}
\end{exer}

\begin{exa}
Let $A\in\mats{3\times 3}$ be given by$$A =
\begin{bmatrix} 1 & 1 & 1 \cr 1 & 1 & 1 \cr 1 & 1 & 1 \cr
\end{bmatrix}.$$
We will demonstrate, using a technique called induction, 
that  $A^n = 3^{n - 1}A$ for $n \in \BBN , n \geq
1$.
 \begin{sol} The assertion is trivial for $n = 1$.
Assume its truth for $n -1$, that is, assume $A^{n - 1} = 3^{n -
2}A$. Observe that
$$A^2 = \begin{bmatrix} 1 & 1 & 1 \cr 1 & 1 & 1 \cr 1 & 1 & 1 \cr
\end{bmatrix}\begin{bmatrix} 1 & 1 & 1 \cr 1 & 1 & 1 \cr 1 & 1 & 1 \cr
\end{bmatrix} = \begin{bmatrix} 3 & 3 & 3 \cr 3 & 3 & 3 \cr 3 & 3 & 3 \cr
\end{bmatrix} = 3A.$$
Now
$$A^n = AA^{n - 1} = A(3^{n - 2}A) = 3^{n - 2}A^2 = 3^{n - 2}3A = 3^{n - 1}A, $$
and so the assertion is proved by induction.
\end{sol}
Do not worry if you do not fully buy into this proof at this point.
Although it may appear we used circular reasoning here, we
in fact did not. All will be explained in the section on induction later.
But we couldn't resist including this example as a foreshadowing
of things to come!
\end{exa}




\begin{thm} 
There is a unique matrix $E\in\mats{n\times n}$ such that 
for every $A\in\mats{n\times n}$, 
$AE = EA = A$. In particular, that
matrix is $E = {\bf I}_n$, the identity matrix. 

\begin{pf}
It is clear that for any $A\in\mats{n\times n}$, $A{\bf I}_n
= {\bf I}_nA = A$. Now because $E$ is an identity, $E{\bf I}_n =
{\bf I}_n$. Because ${\bf I}_n$ is an identity, $E{\bf I}_n =  E$.
Whence
$${\bf
I}_n = E{\bf I}_n =  E,$$demonstrating uniqueness.
\end{pf}
\end{thm}

\begin{rem}
The fact that ${\bf I}_n$ is the multiplicative identity of 
$\mats{n\times n}$
(much like 1 is the multiplicative identity of real numbers) is
the reason we call it the identity matrix.
\end{rem}

\begin{exa}
Let $A = [a_{ij}]\in \mats{n\times n}$ be such that $a_{ij} = 0$
for $i > j$ and $a_{ij} = 1$ if $i \leq j$. Find $A^2$. 
\begin{sol} Let $A^2 = B = [b_{ij}]$. Then
$$b_{ij} = \sum_{k = 1} ^n a_{ik}a_{kj}.$$Observe that the $i$-th
row of $A$ has $i - 1$ $0$'s followed by $n - i + 1$ $1$'s, and
the $j$-th column of $A$ has $j$ $1$'s followed by $n - j$ 0's.
Therefore if $i - 1 > j$, then $b_{ij} = 0$. If $i \leq j + 1$,
then $$b_{ij} = \sum_{k = i} ^{j} a_{ik}a_{kj} = j - i + 1.$$ This
means that $$A^2 = \begin{bmatrix} 1 & 2 & 3 & 4 & \cdots & n - 1
& n \cr 0 & 1 & 2 & 3 & \cdots & n - 2 & n - 1 \cr 0 & 0 & 1 & 2 &
\cdots & n - 3 & n - 2 \cr \vdots & \vdots & \vdots & \vdots &
\cdots &\vdots & \vdots \cr 0 & 0 & 0 & 0 & \cdots & 1 & 2 \cr 0 &
0 & 0 & 0 & \cdots & 0 & 1\cr
\end{bmatrix}.$$
\end{sol}
\end{exa}


\begin{exa}
Let $x$ be a real number, and let  $$m(x) = \begin{bmatrix}1 & 0 &
x\cr -x & 1 & -\dfrac{x^2}{2} \cr 0 &0 &1\cr \end{bmatrix}.  $$If
$a, b$ are real numbers, prove that
\begin{enumerate}
\item $m(a)m(b)=m(a+b)$.
\item $m(a)m(-a)={\bf I}_3$, the $3\times 3$ identity matrix.
\end{enumerate}
\begin{sol}
For the first part, observe that $$\begin{array}{llll}m(a)m(b) & = &
\begin{bmatrix}1 & 0 & a\cr -a & 1 & -\dfrac{a^2}{2} \cr 0 &0 &1\cr
\end{bmatrix}\begin{bmatrix}1 & 0 & b\cr
-b & 1 & -\dfrac{b^2}{2} \cr 0 &0 &1\cr \end{bmatrix} 
 = \begin{bmatrix}1 & 0 & a+b\cr -a-b & 1 &
-\dfrac{a^2}{2}-\dfrac{b^2}{2}+ab \cr 0 &0
&1\cr \end{bmatrix}\\
 &=&  \begin{bmatrix}1 & 0 & a+b\cr -(a+b) & 1 & -\dfrac{(a+b)^2}{2}
\cr 0 &0
&1\cr \end{bmatrix}
=  m(a+b)\end{array}$$ 
For the second part, observe that using
the preceding part of the problem, $$m(a)m(-a)= m(a-a) = m(0)
=\begin{bmatrix}1 & 0 & 0\cr -0 & 1 & -\dfrac{0^2}{2} \cr 0 &0 &1\cr
\end{bmatrix} ={\bf I}_3,$$
giving the result.
\end{sol}
\end{exa}

\begin{exer}
A square matrix $X$ is called {\em idempotent} if $X^2=X$. Prove
that if $AB=A$ and $BA=B$ then $A$ and $B$ are idempotent. 
\vspace*{1.5in}
\begin{answer}
Observe that
$A^2 = (AB)(AB) = A(BA)B = A(B)B = (AB)B = AB =A. $
Similarly,
$B^2 = (BA)(BA) = B(AB)A = B(A)A = (BA)A = BA =B. $
\end{answer}
\end{exer}


\begin{exer}
Prove or disprove: If  $A,B\in\mats{n\times n}$ are
such that $AB = {\bf 0}_n$, then also $BA = {\bf 0}_n.$
Hint: Play around with some simple $2\times 2$ matrices.
\vspace*{1.75in}
\begin{answer}
Disprove! Take $\dis{A = \left[\begin{array}{ll}1 & 0 \\ 0 & 0
\end{array}\right]}$ and $\dis{B = \left[\begin{array}{ll}0 & 1 \\ 0 & 0
\end{array}\right]}$. Then $AB = B$, but
$BA = {\bf 0}_2$.
\end{answer}
\end{exer}

\newpage
\subsection{Trace and Transpose}\label{subsec:matT}
\begin{df}
Let $A = [a_{ij}] \in \mats{n\times n}$. Then the {\em trace}
of $A$, denoted by $\tr{A}$ is the sum of the diagonal elements of
$A$. That is,
$$\tr{A} = \sum _{k = 1} ^n a_{kk}.$$ \index{trace}
\end{df}
\begin{thm}\label{tracethm}
Let $A = [a_{ij}] \in \mats{n\times n}, B = [b_{ij}]\in
\mats{n\times n}$. Then
\begin{equation}  \tr{A + B} = \tr{A} + \tr{B},\end{equation}
\begin{equation}  \tr{AB} = \tr{BA}.\end{equation}

\begin{pf}
The first assertion is trivial. To prove the second, observe that
$AB = (\sum _{k = 1} ^n a_{ik}b_{kj})$ and $BA = (\sum _{k = 1} ^n
b_{ik}a_{kj})$. Then
$$\tr{AB} = \sum _{i = 1} ^n \sum _{k = 1} ^n a_{ik}b_{ki} =
 \sum_{k = 1} ^n \sum _{i = 1} ^n b_{ki}a_{ik} = \tr{BA},
 $$whence the theorem follows.
\end{pf}
\end{thm}

%\begin{exa}
%Let $A\in \mats{n\times n}$. Prove that $A$ can be written as
%the sum of two matrices whose trace is different from $0$.
%\begin{sol} Write $$A = (A - \alpha {\bf I}_n) + \alpha {\bf I}_n.$$ Now,
%$\tr{A - \alpha {\bf I}_n}  = \tr{A} - n\alpha$ and $\tr{\alpha {\bf
%I}_n} = n\alpha$. Thus it suffices to take $\alpha \neq
%\dfrac{\tr{A}}{n}, \alpha \neq 0$. Since $\BBR$ has infinitely many
%elements, we can find such an $\alpha$.
%\end{sol}
%\end{exa}

\begin{exa}
Let $A, B$ be square matrices of size $n>0$. 
Is it possible that $AB - BA = {\bf
I}_n$? Prove or disprove!
\begin{sol} 
This is impossible. Consider taking traces on both sides:
$$0 = \tr{AB}-\tr{BA} = \tr{AB - BA} = \tr{{\bf I}_n} = n$$ 
which is a contradiction, since $n >0$.
\end{sol}
\end{exa}


\begin{exer}
 Consider the matrix $A = \begin{bmatrix} a & b
\cr c & d \cr\end{bmatrix} \in \mats{2\times 2}$. Find
necessary and sufficient conditions on $a, b, c, d$ so that
$\tr{A^2} = (\tr{A})^2$. 
\vspace*{1.5in}
\begin{answer} We have $\tr{A^2} =
\tr{\begin{bmatrix} a & b \cr c & d \cr\end{bmatrix}\begin{bmatrix}
a & b \cr c & d \cr\end{bmatrix}} = \tr{\begin{bmatrix} a^2 + bc &
ab + bd \cr ca + cd & d^2 + cb \cr\end{bmatrix}} = a^2 + d^2 + 2bc
$
and $\left(\tr{\begin{bmatrix} a
& b \cr c & d \cr\end{bmatrix}}\right)^2 = (a+d)^2.  $ Thus $
\tr{A^2} = (\tr{A})^2 \iff a^2 + d^2 + 2bc = (a + d)^2 \iff bc =
ad,
$
is the condition sought.
\end{answer}
\end{exer}

\begin{df} \index{transpose}
The {\em transpose} of a matrix $A = [a_{ij}] \in
\mats{m\times n}$ is the matrix $A^T = B = [b_{ij}] \in
\mats{n\times m}$, where $b_{ij} = a_{ji}$.
\index{matrix!transpose}
\end{df}

\begin{exa}
If $\displaystyle M = \begin{bmatrix}a & b & c \cr  d & e & f \cr g & h & i
\cr\end{bmatrix},$ then
$\displaystyle M^T = \begin{bmatrix}a & d & g \cr  b & e & h \cr c & f & i
\cr\end{bmatrix}.$
\end{exa}

\begin{exer}
Find $N^T$ if  $\displaystyle N = \begin{bmatrix}
1&2&3&4 \cr 
5&6&7&8\cr 
9&10&11&12 \cr \end{bmatrix}.$ 
\vspace*{2in}
\begin{answer}
$\displaystyle N^T = \begin{bmatrix}
1&5&9\cr
2&6&10\cr
3&7&11\cr
4&8&12\cr\end{bmatrix}.$
\end{answer}
\end{exer}

\begin{thm}\label{thm:properties_transpose}
Let $$A = [a_{ij}] \in \mats{m\times n},\  B = [b_{ij}] \in
\mats{m\times n},\  C = [c_{ij}] \in \mats{n\times r},\ \alpha \in \BBR, u\in\BBN .$$
Then
\begin{equation} A^{TT} = A,\end{equation}
\begin{equation} (A + \alpha B)^T =\  A^T + \alpha B^T,\end{equation}
\begin{equation} (AC)^T =\   C^TA^T,\end{equation}
\begin{equation} (A^{u})^T = (A^T)^u.\end{equation}

\begin{pf}
The first two assertions are obvious, and the fourth follows from
the third by using induction. To prove the third put $A^T =
(\alpha_{ij}), \ \alpha _{ij} = a_{ji}$, $C^T = (\gamma_{ij}), \
\gamma _{ij} = c_{ji}$, $AC = (u_{ij})$ and $C^TA^T = (v_{ij})$.
Then
$$u_{ij} = \sum _{k = 1} ^n a_{ik}c_{kj} = \sum _{k = 1} ^n \alpha _{ki}\gamma _{jk} =
\sum _{k = 1} ^n \gamma _{jk}\alpha _{ki} = v_{ji},$$ whence the
theorem follows.
\end{pf}
\end{thm}

\begin{df}\index{matrix!skew-symmetric} \index{matrix!symmetric}
A square matrix $A\in\mats{n\times n}$ is {\em symmetric} if
$A^T = A.$ A matrix $B\in\mats{n\times n}$ is {\em
skew-symmetric} if $B^T = -B.$
\end{df}


\begin{exa} Let $A, B$ be square matrices of the same size, with $A$
symmetric and $B$ skew-symmetric. Prove that the matrix $A^2BA^2$ is
skew-symmetric.
\begin{sol} We have
$$(A^2BA^2)^T = (A^2)^T (B)^T(A^2)^T = A^2(-B)A^2 = -A^2BA^2.$$
\end{sol} 
\end{exa} 

\begin{exer}
Let $A, B$ be square matrices of the same size, with $A$ symmetric
and $B$ skew-symmetric. Prove that the matrix $AB - BA$ is
symmetric. 
\vspace*{1in}
\begin{answer} We have
$(AB - BA)^T = (AB)^T - (BA)^T = B^TA^T - A^TB^T = -BA - A(-B) = AB - BA.$
\end{answer}
\end{exer}
%\begin{thm}
%Let $\BBR$ be a field of characteristic different from $2$.  Then
%any square matrix $A$ can be written as the sum of a symmetric and a
%skew-symmetric matrix.
%\begin{pf} Observe that
%$$(A + A^T)^T = A^T + A^{TT} = A^T + A,$$and so
%$A + A^T$ is symmetric. Also,
%$$(A - A^T)^T = A^T - A^{TT} = -(A - A^T),$$and so $A - A^T$ is
%skew-symmetric. We only need to write $A$ as
%$$A = (2^{-1})(A + A^T) +  (2^{-1})(A - A^T)$$to prove the
%assertion.
%\end{pf}
%\end{thm}

%\begin{exa}
%Find,  with proof, a square matrix $A$ with entries in $\BBZ _2$
%such $A$ is not the sum of a symmetric and and anti-symmetric
%matrix. \end{exa}
%\begin{sol}In $\BBZ_2$ every symmetric matrix is also
%anti-symmetric, since $-x = x$. Thus it is enough to take a
%non-symmetric matrix, for example,  take $A =
%\begin{bmatrix} 0 & 1 \cr 0 &
%0 \cr
%\end{bmatrix}$.
%\end{sol}

